\part{頻度論的統計}
	\newcommand{\xmin}{x_\textrm{min}}
	\newcommand{\xmax}{x_\textrm{max}}
	\newcommand{\ymin}{y_\textrm{min}}
	\newcommand{\ymax}{y_\textrm{max}}
	\chapter{一般の分布}
    \section{\texorpdfstring{$\bm{x}$}{x}の分散共分散行列を\texorpdfstring{$V=\Var{}{\bm{x}}$}{V=Var(x)}とすると\texorpdfstring{$\Var{}{\bm{v}^\top\bm{x}} = \bm{v}^\top V\bm{v}$}{Var(v\string^T x) = v\string^T V v}}
        \begin{proof}
        \quad\par
            $\hat{\bm{x}} = \E{}{\bm{x}}$とすると
            \begin{align*}
                \Var{}{\bm{v}^\top\bm{x}} &= \E{}{\left(\bm{v}^\top\bm{x}-\E{}{\bm{v}^\top\bm{x}}\right)^2} = \E{}{\left(\bm{v}^\top\bm{x}-\bm{v}^\top\hat{\bm{x}}\right)^2} = \E{}{\left(\bm{v}^\top(\bm{x}-\hat{\bm{x}})\right)^2}\\
                &= \E{}{\bm{v}^\top\left(\bm{x}-\hat{\bm{x}}\right)\left(\bm{x}-\hat{\bm{x}}\right)^\top\bm{v}} = \bm{v}^\top\E{}{\left(\bm{x}-\hat{\bm{x}}\right)\left(\bm{x}-\hat{\bm{x}}\right)^\top}\bm{v} = \bm{v}^\top V\bm{v}
            \end{align*}
        \end{proof}
    \section{ベクトル確率変数の線形写像の像の期待値と分散共分散行列}
        \begin{shadebox}
            $X_1,\dotsc,X_n\in\realNumbers$を確率変数、$A\in\realNumbers^{n\times m}$、$\bm{X}\coloneqq[X_1,\dotsc,X_n]^\top$、$\Sigma_X\coloneqq\CovMat{\bm{X}}$とすると次が成り立つ。
            \[\E{}{A^\top\bm{X}} = A^\top\E{}{\bm{X}},\quad\CovMat{A^\top\bm{X}} = A^\top\Sigma_X A\]
            特に$X_1,\dotsc,X_n$が独立で同一の分散$\sigma^2$をもち、かつ$A$が$n$次直交行列であれば
            \[\CovMat{A^\top\bm{X}} = \Sigma_X\]
        \end{shadebox}
        \begin{proof}
            \newcommand{\ai}{{\bm{a}_i}}
            \newcommand{\aj}{{\bm{a}_j}}
            \quad\par
            $\E{}{A^\top\bm{X}} = A^\top\E{}{\bm{X}}$は期待値の線形性より明らか。$\CovMat{A^\top\bm{X}} = A^\top\Sigma_X A$を示す。
            $\bm{Y}\coloneqq A^\top\bm{X}$とし、$A$の第$i$列ベクトルを$\bm{a}_i$とすると
            \begin{align*}
                \CovMat{A^\top\bm{X}}_{ij} &= \CovMat{\bm{Y}}_{ij} = \Cov{Y_i}{Y_j} = \E{}{\left(Y_i-\E{}{Y_i}\right)\left(Y_j-\E{}{Y_j}\right) } \\
                &= \E{}{\left(\ai^\top\bm{X}-\E{}{\ai^\top\bm{X}}\right) \left(\aj^\top\bm{X}-\E{}{\aj^\top\bm{X}}\right)} \\
                &= \E{}{\ai^\top\left(\bm{X}-\E{}{\bm{X}}\right)\aj^\top\left(\bm{X}-\E{}{\bm{X}}\right)} \\
                &= \E{}{\ai^\top\left(\bm{X}-\E{}{\bm{X}}\right)\left(\bm{X}-\E{}{\bm{X}}\right)^\top\aj} \\
                &= \ai^\top\E{}{\left(\bm{X}-\E{}{\bm{X}}\right)\left(\bm{X}-\E{}{\bm{X}}\right)}\aj = \ai^\top\Sigma_X\aj
            \end{align*}
            これより$\CovMat{A^\top\bm{X}} = A^\top\Sigma_X A$。
            \par
            特に$X_1,\dotsc,X_n$が独立で同一の分散$\sigma^2$をもち、かつ$A$が$n$次直交行列であれば
            \[\CovMat{A^\top\bm{X}} = A^\top\Sigma_X A = A^\top\sigma^2I_n A = \sigma^2I_n = \Sigma_X\]
        \end{proof}
    \section{スコア関数}
        $\theta$をパラメータとする確率密度関数$f(x;\theta)$に従うi.i.d確率変数$X_1,\dotsc,X_n \in \Omega$に対して「スコア関数」を次式で定義する。
        \[ V_n(\bm{X},\theta) \coloneqq \partDerivLong{\log f(\bm{X};\theta)}{\theta}{} = \partDerivLong{\log \prod_{i=1}^n f(X_i;\theta)}{\theta}{} = \partDerivLong{\sum_{i=1}^n \log f(X_i;\theta)}{\theta}{} = \sum_{i=1}^n V(X_i,\theta) \]
        \subsection{\texorpdfstring{$\E{\bm{X}}{V(\bm{X},\theta)} = 0$}{E\_X (V(X,Θ)) = 0}}
            \begin{proof}
                \begin{align*}
                    \E{\bm{X}}{V_n(\bm{X},\theta)} &= \integrate{\Omega^n}{}{V_n(\bm{x},\theta) \prod_{i=1}^n f(x_i;\theta)}{}{\bm{x}} = \integrate{\Omega^n}{}{\partDerivLong{\prod_{i=1}^n f(x_i;\theta)}{\theta}{}}{}{\bm{x}} \\
                    &= \partDerivLong{\integrate{\Omega^n}{}{\prod_{i=1}^n f(x_i;\theta)}{}{\bm{x}}}{\theta}{} \quad (\text{微分と積分の順序交換可能性を仮定}) \\
                    &= \partDeriv{1}{\theta}{} = 0
                \end{align*}
            \end{proof}
    \section{Cram\'er-Raoの不等式}
        \begin{shadebox}
            $\theta$をパラメータとする確率密度関数$f(x;\theta)$に従うi.i.d確率変数$X_1,\dotsc,X_n \in \Omega$から計算した$\theta$の不偏推定量$\hat{\theta}(\bm{X})$について次式が成り立つ
            \[ \Var{\bm{X}}{\hat{\theta}(\bm{X})} \geq \left(n\E{\textcolor{red}{X}}{V(\textcolor{red}{X},\theta)^2}\right)^{-1} = \left(n\Var{X}{V(X,\theta)}\right)^{-1} \]
        \end{shadebox}
        \begin{proof}
            \quad\par
            (参考: \url{https://to-kei.net/estimator/cramer-rao/})\\
            \begin{align*}
                \Var{\bm{X}}{\hat{\theta}(\bm{X})}\E{\bm{X}}{V_n(\bm{X},\theta)^2} &= \integrate{\Omega}{}{\left(\hat{\theta}(\bm{x}) - \theta\right)^2 f(\bm{x};\theta)}{}{\bm{x}} \integrate{\Omega}{}{V_n(\bm{x},\theta)^2 f(\bm{x};\theta)}{}{\bm{x}} \\
                &= \integrate{\Omega}{}{\left(\left(\hat{\theta}(\bm{x}) - \theta\right)\sqrt{f(\bm{x};\theta)}\right)^2}{}{\bm{x}} \integrate{\Omega}{}{\left(V_n(\bm{x},\theta)\sqrt{f(\bm{x};\theta)}\right)^2}{}{\bm{x}} \\
                &\geq \integrate{\Omega}{}{\left(\hat{\theta}(\bm{x}) - \theta\right)V_n(\bm{x},\theta)f(\bm{x};\theta)}{}{\bm{x}} \quad (\text{Cauchy-Schwarzの不等式を用いた}) \\
                &= \integrate{\Omega}{}{\hat{\theta}(\bm{x})V_n(\bm{x},\theta)f(\bm{x},\theta)}{}{\bm{x}} \quad \left(\because\; \text{直前の定理より }\E{\bm{X}}{V_n(\bm{X},\theta)} = 0\right) \\
                &= \integrate{\Omega}{}{\hat{\theta}(\bm{x})\left(\partDerivLong{\log f(\bm{x};\theta)}{\theta}{}\right)f(\bm{x};\theta)}{}{\bm{x}} = \integrate{\Omega}{}{\hat{\theta}(\bm{x})\partDerivLong{f(\bm{x};\theta)}{\theta}{}}{}{\bm{x}} \\
                &= \partDerivLong{\integrate{\Omega}{}{\hat{\theta}(\bm{x})f(\bm{x};\theta)}{}{\bm{x}}}{\theta}{} \quad (\text{微分と積分の順序交換可能性を仮定した}) \\
                &= \partDerivLong{\E{\bm{X}}{\hat{\theta}(\bm{X})}}{\theta}{} = \partDeriv{\theta}{\theta}{} = 1
            \end{align*}
            ここで
            \begin{align*}
                \E{\bm{X}}{V_n(\bm{X},\theta)^2} &= \E{\bm{X}}{\left(\sum_{i=1}^n V(X_i,\theta)\right)^2} = \sum_{i=1}^n\sum_{j=1}^n\E{\bm{X}}{V(X_i,\theta)V(X_j,\theta)} \\
                &= \sum_{i=1}^n \E{\bm{X}}{V(X_i,\theta)^2} \quad (\because\; V(X_i,\theta),V(X_j,\theta),\;i\neq j\text{は独立で}\E{\bm{X}}{V(X_i,\theta)} = 0) \\
                &= \sum_{i=1}^n \E{\textcolor{red}{X_i}}{V(X_i,\theta)^2} = n\E{X}{V(X,\theta)^2}
            \end{align*}
            以上より定理の主張が成り立つ。
        \end{proof}
    \section{条件付き確率}
        \subsection{\texorpdfstring{$\Pr{A}=\sum_k{\Pr{A,B_k}} = \sum_k{\cPr{A}{B_k}\Pr{B_k}}$}{Pr(A) = Σ\_k Pr(A, Bk) = Σ\_k Pr(A|Bk)Pr(Bk)}}
        \label{tech:確率の分解}
        \subsection{\texorpdfstring{$\cPr{A,B}{C}=\cPr{A}{B,C}\cPr{B}{C}$}{Pr(A,B|C) = Pr(A|B,C)Pr(B|C)} (スライディング)}
            \begin{proof}
                \[\cPr{A,B}{C} = \frac{\Pr{A,B,C}}{\Pr{C}} = \frac{\cPr{A}{B,C}\Pr{B,C}}{\Pr{C}} = \cPr{A}{B,C}\cPr{B}{C}\]
            \end{proof}
        \subsection{\texorpdfstring{$\cPr{A}{C} = \sum_k{\cPr{A}{B_k,C}\cPr{B_k}{C}}$}{Pr(A|C) = Σ\_k Pr(A|Bk,C)Pr(Bk|C)}}
            \begin{proof}
                \begin{align*}
                    \cPr{A}{C} &= \frac{\Pr{A,C}}{\Pr{C}} = \frac{\sum_k{\Pr{A,B_k,C}}}{\Pr{C}} \quad(\because\;\ref{tech:確率の分解}) \\
                    &= \frac{\sum_k{\cPr{A}{B_k,C}\Pr{B_k,C}}}{\Pr{C}} = \sum_k{\cPr{A}{B_k,C}\cPr{B_k}{C}}
                \end{align*}
            \end{proof}
    \section{条件付き確率密度関数}
        連続型2変数$X,Y$に関して、条件付き確率密度関数を次のように定義する。
        \[p_Y(y|X=x) \coloneqq \frac{p(x,y)}{\integrate{\ymin}{\ymax}{p(x,y)}{}{y}}\]
        ($p(x,y)$は$X,Y$の同時確率密度関数)
    \section{条件付き期待値}
        2変数$X,Y$に関して、条件付き期待値を次のように定義する。
        \begin{align*}
            \textrm{(離散)} &: \E{Y}{Y|X=x} \coloneqq \sum_y y\Pr{Y=y|X=x} = \sum_y y\frac{\Pr{Y=y,X=x}}{\Pr{X=x}} \\
            \textrm{(連続)} &: \E{Y}{Y|X=x} \coloneqq \integrate{\ymin}{\ymax}{yp_Y(y|X=x)}{}{y}
        \end{align*}
        \subsection{\texorpdfstring{$\E{X,Y}{\E{Y}{Y|X}} = \E{X,Y}{Y}$}{E\_(X,Y) (E\_Y (Y|X)) = E\_(X,Y) (Y)}}
            \begin{shadebox}
                条件付き期待値について次の性質が成り立つ。
                \[\E{X,Y}{\E{Y}{Y|X}} = \E{X,Y}{Y}\]
            \end{shadebox}
            \begin{proof}
                \quad\\
                (離散型)
                \begin{align*}
                    \E{X,Y}{\E{Y}{Y|X}} &= \sum_x\sum_y\E{Y}{Y|X=x}\Pr{X=x,Y=y} \\
                    &= \sum_x\E{Y}{Y|X=x}\sum_y\Pr{X=x,Y=y} \quad (\because\;\E{Y}{Y|X=x}\textrm{は}y\textrm{に無関係}) \\
                    &= \sum_x\left(\sum_y{\frac{\Pr{X=x,Y=y}}{\Pr{X=x}}}\right)\Pr{X=x} \\
                    &= \sum_x\sum_y y\Pr{X=x,Y=y} = \E{X,Y}{Y}
                \end{align*}
                (連続型)\\
                $X,Y$の同時確率密度関数を$p(x,y)$とし、$X$の周辺確率密度関数を$p_X(x)$とする。
                \begin{align*}
                    \E{X,Y}{\E{Y}{Y|X}} &= \integrate{\xmin}{\xmax}{\integrate{\ymin}{\ymax}{\E{Y}{Y|X=x}p(x,y)}{}{y}}{}{x} \\
                    &= \integrate{\xmin}{\xmax}{\E{Y}{Y|X=x}\integrate{\ymin}{\ymax}{p(x,y)}{}{y}}{}{x} \quad (\because\;\E{Y}{Y|X=x}\textrm{は}y\textrm{に無関係}) \\
                    &= \integrate{\xmin}{\xmax}{\left(\integrate{\ymin}{\ymax}{y\frac{p(x,y)}{p_X(x)}}{}{y}\right)p_X(x)}{}{x} = \integrate{\xmin}{\xmax}{\integrate{\ymin}{\ymax}{yp(x,y)}{}{y}}{}{x} = \E{X,Y}{Y}
                \end{align*}
            \end{proof}
    % \begin{comment}
    % \section{指示関数}
    %     \subsection{同時事象の指示関数の期待値}
    %         \begin{shadebox}
    %             同時事象の指示関数の期待値に関して次が成り立つ。
    %             \begin{align*}
    %                 \E{}{\ind{Y \satisfy C_Y,\;X \satisfy C_X}} &= \E{X}{\ind{X\satisfy C_X}\E{Y}{\cind{Y\satisfy C_Y}{X\satisfy C_X}}} \\
    %                 &= \E{X}{\ind{X\satisfy C_X}\cPr{Y\satisfy C_Y}{X\satisfy C_X}}
    %             \end{align*}
    %         \end{shadebox}
    %         \begin{proof}
    %             \begin{align*}
    %                 \textrm{左辺} &= \E{}{\cind{Y\satisfy C_Y}{X\satisfy C_X}\ind{X\satisfy C_X}} \\
    %                 &= \sum_x{\sum_y{\Pr{X=x,Y=y}\cind{y\satisfy C_Y}{x\satisfy C_X}\ind{x\satisfy C_X}}} \\
    %                 &= \sum_x{\sum_y{\cPr{Y=y}{X=x}\Pr{X=x}\cind{y\satisfy C_Y}{x\satisfy C_X}\ind{x\satisfy C_X}}} \\
    %                 &= \sum_x{\Pr{X=x}\ind{x\satisfy C_X}\sum_y{\cPr{Y=y}{X=x}\cind{y\satisfy C_Y}{x\satisfy C_X}}} \\
    %                 &= \sum_x{\Pr{X=x}\ind{x\satisfy C_X}\E{Y}{\cind{Y\satisfy C_Y}{x\satisfy C_X}}} \\
    %                 &= \E{X}{\ind{X\satisfy C_X}\E{Y}{\cind{Y\satisfy C_Y}{X\satisfy C_X}}} \\
    %                 &= \E{X}{\ind{X\satisfy C_X}\cPr{Y\satisfy C_Y}{X\satisfy C_X}}
    %             \end{align*}
    %         \end{proof}
    % \end{comment}
    \section{確率不等式}
        \subsection{Markovの不等式}
            \begin{shadebox}
                $(X,\mathcal{M},\mu)$を測度空間、$f:X\to[-\infty,\infty]$を可測関数、$a>0$とするとき、次が成り立つ。
                \[\displaystyle \mu(\{x\in X\,|\,|f(x)|\ge a\})\le\frac{1}{a}\int_{X}|f(x)|\mu(dx)\]
            \end{shadebox}
            \begin{proof}
                \quad\par
                (引用 : \url{http://mathcommunication.hatenablog.com/entry/2017/06/27/001722})
                \par
                \begin{equation*}
                    \begin{split}
                        \int_{X}|f(x)|\mu(dx) & =\int_{|f(x)|\ge a} |f(x)|\mu(dx) + \int_{|f(x)|<a}|f(x)|\mu(dx) \\
                        & \ge\int_{|f(x)|\ge a} |f(x)|\mu(dx) \\
                        & \ge a\,\mu(\{x\in X\,|\,|f(x)|\ge a\})
                    \end{split}
                \end{equation*}
            \end{proof}
        \subsection{Chebyshevの不等式}
            \begin{shadebox}
                任意の連続型確変数$X$について次が成り立つ。
                \[ \Pr{|X-\mu|\geq a\sigma} \leq 1/a^2 \]
                但し$a>0, \mu = \E{}{X}, \sigma^2 = \Var{}{}{X}$
            \end{shadebox}
            \begin{proof}
                \quad\par
                $X$の確率密度関数を$f(x)$とすると
                \begin{align*}
                    \Pr{|X-\mu| \geq a\sigma} &= \left(\int_{-\infty}^{\mu-a\sigma} + \int_{\mu+a\sigma}^\infty\right) \leq \left(\int_{-\infty}^{\mu-a\sigma} + \int_{\mu+a\sigma}^\infty\right)\frac{(x-\mu)^2}{a^2\sigma^2}f(x)\mathrm{d}x \\
                    &\left(\because\;\frac{(x-\mu)^2}{a^2\sigma^2} \geq 1, x\in(-\infty,\mu-a\sigma]\cup[\mu+a\sigma,\infty)\right) \\
                    &\leq \integrate{-\infty}{\infty}{\frac{(x-\mu)^2}{a^2\sigma^2}f(x)}{}{x} = 1/a^2
                \end{align*}
            \end{proof}
        \subsection{H\"{o}ffdingの補題}
            \begin{shadebox}
                $(\Omega,\mathcal{F},P)$を確率空間とし、$X$を$E_{P}[X]=0$、$a\leq X\leq b$ (a.e.)を満たす確率変数とする。このとき任意の$t>0$に対して次の不等式が成り立つ。
                \[E_{P}[\exp(tX)]\le\exp\left(\frac{1}{8}t^{2}(b-a)^{2}\right)\]
            \end{shadebox}
            \begin{proof}
                \quad\par
                (参考 : \url{http://mathcommunication.hatenablog.com/entry/2017/06/27/001722})
                \par
                $E_{P}[X]=0$より$a<0,b>0$であることに注意。指数関数の凸性より、
                \[\exp(tX)\le\frac{X-a}{b-a}\exp(tb)+\frac{b-X}{b-a}\exp(ta),\quad(\mathrm{a.e.})\]
                であるから
                \begin{equation}
                    E_{P}[\exp(tX)]\le\frac{-a}{b-a}\exp(tb)+\frac{b}{b-a}\exp(ta)
                \end{equation}
                今、
                \[\displaystyle p\coloneqq\frac{-a}{b-a},\quad u\coloneqq t(b-a)\]
                とおくと$\displaystyle p\coloneqq\frac{-a}{b-a},\quad u\coloneqq t(b-a)$であり、
                \begin{align*}
                    (1) &= p\exp(tb)+(1-p)\exp(ta)\\
                    &= p\exp((1-p)u)+(1-p)\exp(-pu)
                \end{align*}
                \begin{screen}
                    $\because\;\;1-p=\frac{b}{b-a}$より$tb = t(b-a)(1-p) = (1-p)u$ また、$ta = -tp(b-a) = -pu$
                \end{screen}
                さらに、$\phi(u) \coloneqq \log{\left(\textrm{式}(1)\right)}$とおくと
                \begin{align*}
                    \phi(u) &= \log\left[p\exp( (1-p)u)+(1-p)\exp(-pu)\right]\\
                    &= \log\left[\exp( (1-p)u)\cdot(p+(1-p)\exp(-u))\right]\\
                    &= (1-p)u+\log(p+(1-p)\exp(-u))
                \end{align*}
                \begin{gather*}
                    \phi'(u)=1-p+\frac{-(1-p)\exp(-u)}{p+(1-p)\exp(-u)},\\
                    \phi''(u)=\frac{p(1-p)\exp(-u)}{(p+(1-p)\exp(-u))^{2}}
                \end{gather*}
                ここでTaylorの定理より、ある$u_0\in(0,u)$が存在して
                \[\phi(u)=\phi(0)+\phi'(0)u+\frac{1}{2}\phi''(u_{0})u^{2}=\frac{1}{2}\phi''(u_{0})u^{2}\]
                であるが、相加・相乗平均の関係より
                \[\phi''(u_{0})=\frac{p\cdot(1-p)\exp(-u_{0})}{(p+(1-p)\exp(-u_{0}))^{2}}\le\frac{1}{4}\]
                であるので
                \[\phi(u)\le\frac{1}{8}u^{2}\]
                よって
                \[E_{P}[\exp(tX)]\le\exp(\phi(u) )\le\exp\left(\frac{1}{8}u^{2}\right)=\exp\left(\frac{1}{8}t^{2}(b-a)^{2}\right)\]
            \end{proof}
        \subsection{H\"{o}ffdingの不等式}
            \begin{shadebox}
                $(\Omega,\mathcal{F},P)$を確率空間とし、$X_{1},\ldots,X_{n}$を独立同分布な確率変数で$a_{i}\le X_{i} \le b_{i}\ (\mathrm{a.e.})$を満たすとする。このとき任意の$\epsilon>0$に対して次の不等式が成り立つ。
                \begin{gather*}
                    P(S-E_{P}[S]\ge\varepsilon)\le \exp\left(-\frac{2\varepsilon^{2}}{\sum_{i=1}^{n}(b_{i}-a_{i})^{2}}\right),\\
                    P(S-E_{P}[S]\le-\varepsilon)\le\exp\left(-\frac{2\varepsilon^{2}}{\sum_{i=1}^{n}(b_{i}-a_{i})^{2}}\right)
                \end{gather*}
            \end{shadebox}
            \begin{proof}
                \quad\par
                (引用 : \url{http://mathcommunication.hatenablog.com/entry/2017/06/27/001722})
                \par
                任意の$t>0$に対して
                \begin{align*}
                    P(S-E_{P}[S]\ge\varepsilon) &= P(\exp(t(S-E_{P}[S]))\ge\exp(t\varepsilon))\\
                    &\le \exp(-t\varepsilon)E_{P}[\exp(t(S-E_{P}[S]))] & \text{（マルコフの不等式）}\\
                    &= \exp(-t\varepsilon)\prod_{i=1}^{n}E_{P}[\exp(t(X_{i}-E_{P}[X_{i}]))] & (\text{各}X_i\text{は独立})\\
                    &\leq \exp(-t\varepsilon)\prod_{i=1}^{n}\exp\left(\frac{t^2}{8}(b_i-a_i)^2\right) & \text{(ヘフディングの補題)}\\
                    &= \exp(-t\varepsilon)\exp\left(\frac{t^2}{8}\sum_{i=1}^{n}(b_{i}-a_{i})^{2}\right)\\
                    &= \exp\left(\frac{t^2}{8}\sum_{i=1}^{n}(b_{i}-a_{i})^{2} - t\varepsilon\right)
                \end{align*}
                が成り立つ。
                \[\frac{1}{8}t^{2}\sum_{i=1}^{n} (b_{i}-a_{i})^{2}-t\varepsilon = \frac{1}{8}\left(\sum_{i=1}^{n}(b_{i}-a_{i})^{2}\right)\left(t-\frac{4\varepsilon}{\sum_{i=1}^{n}(b_{i}-a_{i})^{2}}\right)^{2}-\frac{2\varepsilon^{2}}{\sum_{i=1}^{n}(b_{i}-a_{i})^{2}}\]
                であり、第1項$\ge0$だから
                \[P(S-E_{P}[S]\ge\varepsilon)\le\exp\left(-\frac{2\varepsilon^{2}}{\sum_{i=1}^{n}(b_{i}-a_{i})^{2}}\right)\]
                同様にして
                \[\displaystyle P(S-E_{P}[S]\le-\varepsilon)\le\exp\left(-\frac{2\varepsilon^{2}}{\sum_{i=1}^{n}(b_{i}-a_{i})^{2}}\right)\]
            \end{proof}
    \section{中心極限定理}
        \renewcommand{\A}{\frac{\sqrt{n}(\hat{\mu}_n-\mu)}{\sigma}}
        \newcommand{\SUM}{\sum_{j=1}^n}
        \newcommand{\PI}{\prod_{j=1}^n}
        \begin{shadebox}
            標準化された標本平均$\A$の分布は標準正規分布に弱収束する。すなわち
            \[\lim_{n\rightarrow\infty}\Pr{\A\leq x} = \Phi(x)\]
            但し$\Phi(x)$は標準正規分布の累積分布関数である。
        \end{shadebox}
        \begin{proof}
            \quad\par
            (引用 : \url{https://ja.wikipedia.org/wiki/中心極限定理})
            \par
            \[Z_n \coloneqq \A = \frac{\SUM{X_j}-n\mu}{\sqrt{n}\sigma} = \SUM{Y_j/sqrt{n}}\quad\textrm{where }Y_j \coloneqq \frac{X_j-\mu}{\sigma}\]
            とおく。$Y_j$の平均は0分散は1である。
            \par
            $Z_n$の特性関数$\phi_{Z_n}(t)$は
            \begin{align*}
                \phi_{Z_n}(t) &= \E{}{e^{itZ_n}} = \E{}{e^{it\SUM{Y_j/\sqrt{n}}}} = \E{}{\PI{e^{itY_j/\sqrt{n}}}} = \PI{\E{}{e^{itY_j/\sqrt{n}}}}\quad(\because\textrm{各}Y_j\textrm{は独立})\\
                &= \PI{\phi_{Y_j}(t/\sqrt{n})} = \left[\phi_{Y_1}(t/\sqrt{n})\right]^n\quad(\because\textrm{各}Y_j\textrm{は全部同等})
            \end{align*}
            $\phi_{Y_1}(t)$のTaylor展開を考える。
            \begin{align*}
                \phi_{Y_1}(0) &= 1\\
                {\phi_{Y_1}}'(0) &= \left.\E{}{iY_1e^{itY_i}}\right|_{t=0} = 0\\
                {\phi_{Y_1}}''(0) &= \left.\E{}{-{Y_1}^2e^{itY_i}}\right|_{t=0} = -\E{}{{Y_1}^2} = -1
            \end{align*}
            これより、
            \[\phi_{Y_1}(t/\sqrt{n}) = 1 - \frac{1}{2}\frac{t^2}{n} + c\frac{1}{3!}\left(\frac{t}{\sqrt{n}}\right)^3 + o\left(\left(\frac{t}{\sqrt{n}}\right)^3\right)\quad(t\rightarrow0)\]
            よって
            \[\phi_{Z_n}(t) = \left[1 - \frac{1}{n}\frac{t^2}{2} + c\frac{1}{3!}\left(\frac{t}{\sqrt{n}}\right)^3 + o\left(\left(\frac{t}{\sqrt{n}}\right)^3\right)\right]^n \rightarrow e^{-t^2/2}\quad\textrm{as }n\rightarrow\infty\]
            最右辺は標準正規分布の特性関数であり、特性関数と確率密度関数は1対1で対応するから$Z_n = \A$は$n\rightarrow\infty$で標準正規分布に従う。
        \end{proof}
	\chapter{多項分布}
    以下 $(X_1,X_2,\dots,X_n) \sim B_N(n,p_1,p_2,\dots,p_n)$ とする。
    \section{周辺分布は2項分布}
        \begin{shadebox}
            \[P(X_h = x_h) = \frac{n!}{x_h!(n-x_h)!}p_h^{x_h}(1-p_h)^{x_h}\]
        \end{shadebox}
        \begin{proof}
            \begin{align*}
                P(X=x_h) &= \sum_{x_1=0}^{n-x_h}\sum_{x_2=0}^{n-x_h-x_1}\cdots\sum_{x_{h-1}=0}^{n-x_h-\sum_{j=1}^{h-2}{x_j}}\sum_{x_{h+1}=0}^{n-\sum_{j=1}^{h}{x_j}}\cdots\sum_{x_n=0}^{n-\sum_{j=1}^{n-1}{x_j}}{P(X_1=x_1,\dots,X_n=x_n)}\\
                &= \sum_{x_1=0}^{n-x_h}\sum_{x_2=0}^{n-x_h-x_1}\cdots\sum_{x_{h-1}=0}^{n-x_h-\sum_{j=1}^{h-2}{x_j}}\sum_{x_{h+1}=0}^{n-\sum_{j=1}^{h}{x_j}}\cdots\sum_{x_n=0}^{n-\sum_{j=1}^{n-1}{x_j}} \\
                &\quad \frac{n!}{\left(\prod_{i=1}^{n}{x_i!}\right)\left(n-\sum_{i=1}^{n}{x_i}\right)!}\left(\prod_{i=1}^{n}{{p_i}^{x_i}}\right)\left(1-\sum_{i=1}^{n}{p_i}\right)^{n-\sum_{i=1}^{n}{x_i}}\\
                &= \frac{n!}{x_h!(n-x_h)!}{p_h}^{x_h}\cdots\sum_{x_{h-1}=0}^{n-x_h-\sum_{j=1}^{h-2}{x_j}}{\frac{\left(n-x_h-\sum_{j=1}^{h-2}{x_j}\right)!}{x_{h-1}!\left(n-x_h-\sum_{j=1}^{h-2}{x_j}-x_{h-1}\right)!}{p_{h-1}}^{x_{h-1}}}\\
                &{} \sum_{x_{h+1}=0}^{n-\sum_{j=1}^{h}{x_j}}{\frac{\left(n-\sum_{j=1}^{h}{x_j}\right)!}{x_{h+1}!\left[\left(n-\sum_{j=1}^{h}{x_j}\right)-x_{h+1}\right]!}{p_{h+1}}^{x_{h+1}}}\cdots\\
                &{} \sum_{x_n=0}^{n-\sum_{j=1}^{n-1}{x_j}}{\frac{\left(n-\sum_{j=1}^{n-1}{x_j}\right)!}{x_n!\left[\left(n-\sum_{j=1}^{n-1}{x_j}\right)-x_n\right]!}{p_n}^{x_n}}\left[\left(1-\sum_{j=1}^{n-1}{p_j}\right)-p_n\right]^{\left[\left(n-\sum_{j=1}^{n-1}{x_j}\right)-x_n\right]}
            \end{align*}
            式の末尾の$\Sigma$は
            \[\left[\left(1-\sum_{j=1}^{n-1}{p_j}\right) - p_n + p_n\right]^{n-\sum_{j=1}^{n-1}{x_j}} = \left[1-\sum_{j=1}^{n-1}{p_j}\right]^{n-\sum_{j=1}^{n-1}{x_j}}\]
            と等しい。同様のことを繰り返すと次の結論に至る。
            \[P(X_h = x_h) = \frac{n!}{x_h!(n-x_h)!}p_h^{x_h}(1-p_h)^{x_h}\]
        \end{proof}
    \section{共分散}
        \begin{shadebox}
            \[\Cov{X_i}{X_j} = -np_ip_j\]
        \end{shadebox}
        \begin{proof}
            \quad\par
            一般に$\Cov{X}{Y} = \E{}{XY}  -\E{}{X}\E{}{Y}$であって、$\E{}{X_i} = np_i, \E{}{X_j} = np_j$は既知であるから$\E{}{X_iX_j}$を計算すれば良い。定理1より
            \[P(X_i=x_i,X_j=x_j) = \frac{n!}{x_i!x_j!(n-x_i-x_j)!}{p_i}^{x_i}{p_j}^{x_j}(1-p_i-p_j)^{n-x_i-x_j}\]
            であることを利用する。$X_i+X_j \leq n$であるから一方が0或いは$n$であれば積が0になることに注意して
            \begin{align*}
                \E{}{X_iX_j} &= \sum_{x_i=\textcolor{red}{1}}^{\textcolor{red}{n-1}}{\sum_{x_j=\textcolor{red}{1}}^{\textcolor{red}{n-x_i}}{x_ix_jP(X_i=x_i,X_j=x_j)}}\\
                &= \sum_{x_i=1}^{n-1}{\sum_{x_j=1}^{n-x_i}{x_ix_j\frac{n!}{x_i!x_j!(n-x_i-x_j)!}{p_i}^{x_i}{p_j}^{x_j}(1-p_i-p_j)^{n-x_i-x_j}}}\\
                &= \sum_{x_i=1}^{n-1}{\frac{n!}{(x_i-1)!(n-x_i)!}{p_i}^{x_i}\sum_{x_j=1}^{n-x_i}{\frac{(n-x_i)!}{(x_j-1)!(n-x_i-x_j)!}{p_j}^{x_j}(1-p_i-p_j)^{n-x_i-x_j}}}\\
                &= \sum_{x_i=1}^{n-1} \frac{n!}{(x_i-1)!(n-x_i)!}{p_i}^{x_i}(n-x_i)p_j \\
                &\quad \sum_{{x_j}'=0}^{n-x_i-1}{\frac{(n-x_i-1)!}{{x_j}'!(n-x_i-1-{x_j}')!}{p_j}^{{x_j}'}(1-p_i-p_j)^{n-x_i-1-{x_j}'}}\\
                &= \sum_{x_i=1}^{n-1}{\frac{n!}{(x_i-1)!(n-x_i)!}{p_i}^{x_i}(n-x_i)p_j(1-p_i)^{n-x_i-1}}\\
                &= p_ip_jn(n-1)\sum_{{x_i}'=0}^{n-2}{\frac{(n-2)!}{{x_i}'!(n-2-{x_i}')!}{p_i}^{{x_i}'}(1-p_i)^{n-2-{x_i}'}}\\
                &= p_ip_jn(n-1)
            \end{align*}
            よって
            \[\Cov{X_i}{X_j} = \E{}{X_iX_j} - \E{}{X_i}\E{}{x_j} = p_ip_jn(n-1) - np_inp_j =  -np_ip_j\]
        \end{proof}
	\chapter{幾何分布}
    \section{無記憶性からの導出}
    \begin{proof}
        \quad\par
        確率変数$X$を、ベルヌーイ試行を繰り返して初めて成功するまでに要する試行回数とする。
        無記憶であるとはすなわち「試行開始から$n$回失敗したという事象の下、それからさらに$k\in\naturalNumbers$回失敗する条件付き確率」が「試行開始から連続で$k$回失敗する確率」に等しいということである。
        特に$k=1$としたときの条件から次の式が成り立つ。
        \begin{align*}
            \frac{P(X>n+1)}{P(X>n)} &= P(X>1) \\
            P(X>n+1) &= P(X>1)P(X>n) \\
            P(X>n) - P(X>n+1) &= P(X>n) - P(X>1)P(X>n) \\
            P(X=n+1) &= (1-P(X>1))P(X>n) = P(X=1)(1-P(X\leq n))
        \end{align*}
        これより
        \begin{align*}
            P(X=2) &= P(X=1)(1 - P(X=1)) = p(1-p) \quad (p \coloneqq P(X=1)) \\
            P(X=3) &= P(X=1)(1 - P(X=1) - P(X=2)) = p[1-p-p(1-p)] = p(1-p)^2 \\
            P(X=4) &= P(X=1)(1 - P(X=1) - P(X=2) - P(X=3)) \\
            &= p[1-p-p(1-p)-p(1-p)^2] = p(1-p)^3 \\
            &\vdots \\
            P(X=n) &= p(1-p)^{n-1}
        \end{align*}
    \end{proof}
	\chapter{指数分布}
    \section{定義}
        $\mu>0$とする。
        確率密度関数$f(x)$が次式となる分布を、尺度母数$\mu$の指数分布という。
        \[
            f(x) =
            \begin{cases}
                \frac{1}{\mu}e^\frac{-x}{\mu} & (x\geq 0) \\
                0 & (x<0)
            \end{cases}
        \]
        確率変数$X$がこの分布に従うことを$X \sim \ExpDist{\mu}$と表すことにする。
    \section{解釈}
        単位時間当たりの生起回数の期待値が$\lambda$であるような事象が時間$x>0$の間に生起する確率の累積分布。
    \section{無記憶性からの導出}
        \begin{proof}
            \quad\par
            「観測開始から時間$s$の間イベントが起こらなかったという事象の下、それから時間$\Delta t$以内にイベントが起こる条件付き確率」が「観測開始($t=0$)から時間$\Delta t$以内にイベントが起こる確率」に等しいという条件より次式が成り立つ。
            \begin{align*}
                \frac{\integrate{s}{s+\Delta t}{f(x)}{}{x}}{\integrate{s}{\infty}{f(x)}{}{x}} &= \integrate{0}{\Delta t}{f(x)}{}{x} \\
                \therefore\; \integrate{s}{s+\Delta t}{f(x)}{}{x} &= \integrate{0}{\Delta t}{f(x)}{}{x} \integrate{s}{\infty}{f(x)}{}{x} \\
                \therefore\; f(s+\Delta t) - f(s) &= -f(s)\integrate{0}{\Delta t}{f(x)}{}{x} \quad (\text{両辺を}s\text{で微分した}) \\
                \therefore\; f'(s) &= -f(0)f(s) \quad (\text{両辺を}\Delta t\text{で割った後}\Delta t \to +0\text{とした})
            \end{align*}
            この微分方程式を解いて規格化すると$f(x) = f(0)\exp(-f(0)x)\;(f(0)>0)$となる。$\lambda = f(0)$とすれば$f(x) = \lambda e^{-\lambda x}$となる。
        \end{proof}
    \section{特性関数}
        \label{指数分布の特性関数}
        尺度母数$\mu$の指数分布の特性関数は次式である。
        \begin{align*}
            \E{}{e^{itX}} &= \integrate{0}{\infty}{\frac{1}{\mu}e^\frac{-x}{\mu}e^{itx}}{}{x} = \frac{1}{\mu}\integrate{0}{\infty}{e^{\left(\frac{-1}{\mu} + it\right)x}}{}{x} = \frac{1}{\mu}\frac{1}{\frac{-1}{\mu} + it}\left[e^{\left(\frac{-1}{\mu} + it\right)x}\right]_0^\infty = \frac{1}{1-i\mu t}
        \end{align*}
	\chapter{Erlang分布}
		\section{定義}
			$k\in\naturalNumbers,\;\mu>0$とする。
			確率密度関数$f(x)$が次式となる分布を、形状母数$k$,尺度母数$\mu$のErlang分布という。
			\[
				f(x) =
				\begin{cases}
					\frac{x^{k-1}}{(k-1)!\mu^k}e^\frac{-x}{\mu} & (x\geq 0) \\
					0 & (x<0)
				\end{cases}
			\]
			確率変数$X$がこの分布に従うことを$X \sim \ErlangDist{k}{\mu}$と表すことにする。
		\section{解釈}
			単位時間当たりの生起回数の期待値が$\lambda = 1/\mu$である事象が、時間$x$の間に$k$以上回数生起する(つまり、$k$回目の事象の発生までの待ち時間が$x$である)確率の累積分布。
		\section{指数分布との関係: \texorpdfstring{$X_1,\dotsc,X_k (\text{独立}) \sim \ExpDist{\mu} \Rightarrow X_1+\dotsb+X_k \sim \ErlangDist{k}{\mu}$}{X1, ..., Xk が互いに独立で母数μの指数分布に従う ⇒ X1+...+Xk は母数μのk次のErlang分布に従う。}}
			\label{指数分布とErlang分布の関係}
			\begin{proof} (直接的な証明)
				\par
				数学的帰納法で示す。
				尺度母数$\mu$の指数分布の確率密度関数を$f$とし、$\sum_{l=1}^k X_l$が従う分布の確率密度関数を$f_k$とする。
				まず$k=1$のときに成り立つことは簡単に確かめられる。
				$k=l\in\naturalNumbers$まで成り立つと仮定して$k=l+1$のときに成り立つことを示す。
				\begin{align*}
					f_{l+1}(x) &= \integrate{0}{x}{f_l(y)f(x-y)}{}{y} = \integrate{0}{x}{\frac{y^{l-1}}{(l-1)!\mu^l}e^\frac{-y}{\mu}\frac{1}{\mu}e^\frac{-(x-y)}{\mu}}{}{y} \\
					&= \frac{1}{(l-1)!\mu^{l+1}}e^\frac{-x}{\mu}\integrate{0}{x}{y^{l-1}}{}{y} = \frac{x^l}{l!\mu^{l+1}}e^\frac{-x}{\mu}
				\end{align*}
			\end{proof}
			\begin{proof} (特性関数を用いた証明)
				\par
				$X_1+\dotsb+ X_k$が従う分布の特性関数は次式である($X_1,\dotsc,X_k$が独立であることに注意)。
				\[ \E{}{e^{it(X_1+\dotsb+X_k)}} = \prod_{l=1}^k\E{}{e^{itX_l}} = \E{}{e^{itX_1}}^k = \frac{1}{(1-i\mu t)^k} \quad (\text{\ref{指数分布の特性関数}}を用いた) \]
				一方、$\ErlangDist{k}{\mu}$の特性関数を計算すると次式になる。
				\begin{align*}
					&\quad \integrate{0}{\infty}{\frac{x^{k-1}}{(k-1)!\mu^k}e^\frac{-x}{\mu}e^{itx}}{}{x} = \frac{1}{(k-1)!\mu^k}\integrate{0}{\infty}{x^{k-1}e^{\left(\frac{-x}{\mu} + it\right)x}}{}{x} \\
					&= \frac{1}{(k-1)!\mu^k}\left[(-1)^{k+1}(k-1)!\frac{1}{\left(\frac{-1}{\mu}+it\right)^k}e^{\left(\frac{-x}{\mu} + it\right)x}\right]_0^\infty \quad (瞬間部分積分を用いた) \\
					&= \frac{1}{(1-i\mu t)^k}
				\end{align*}
				両者は一致する。
			\end{proof}
	\chapter{Poisson分布}
		\section{定義}
			$k\in\naturalNumbers\cup\{0\},\;\mu>0,\;\lambda=1/\mu$とする。
			確率質量関数$P(X=k)$が次式となる分布を、母数$\lambda$のPoisson分布という。
			\[
				P(X=k) =
				\begin{cases}
					\frac{\lambda^k}{k!}e^{-\lambda} & (k = 0,1,2,\dotsc) \\
					0 & (\text{otherwise})
				\end{cases}
			\]
			確率変数$X$がこの分布に従うことを$X \sim \PoissonDist{k}{\mu}$と表すことにする。
		\section{解釈}
			単位時間当たりの生起回数の期待値が$\lambda = 1/\mu$である事象が単位時間内に$k$回生起する確率。
			\subsection{指数分布からの導出}
				指数分布とErlang分布の関係を用いる。\ref{指数分布とErlang分布の関係}において$x=1$(= 単位時間)とすればよい。
		\section{Erlang分布との関係}
			単位時間当たりの生起回数の期待値が$\lambda = 1/\mu$である事象が単位時間内に$k$回以上生起する確率をPoisson分布，Erlang分布の両方の視点から求めてみる。
			Poisson分布の視点で求めると次式となる。
			\[
				\sum_{l=k}^\infty \frac{\lambda^l}{l!}e^{-\lambda}
			\]
			一方でErlang分布の視点で求めると次式となり、当然だが上式と一致する。
			\begin{align*}
				&\integrate{0}{1}{\frac{1}{(k-1)!\mu^k}x^{k-1}e^{-x/\mu}}{}{x} = \frac{1}{(k-1)!\mu^k} \integrate{0}{1}{x^{k-1}e^{-x/\mu}}{}{x} \\
				&\quad = \frac{1}{(k-1)!\mu^k}\left[ -\sum_{l=0}^{k-1}\perm{k-1}{l}\mu^{l+1}x^{k-1-l}e^{-x/\mu} \right]_0^1 = \left[ -\sum_{l=0}^{k-1} \frac{1}{(k-1-l)!\mu^{k-1-l}}x^{k-1-l}e^{-x/\mu} \right]_0^1 \\
				&= \left[ -\sum_{l=0}^{k-1} \frac{x^l}{l!\mu^l}e^{-x/\mu} \right]_0^1 = 1 - \sum_{l=0}^{k-1} \frac{1}{l!\mu^l}e^{-1/\mu} = \sum_{l=0}^\infty \frac{\lambda^l}{l!}e^{-\lambda} - \sum_{l=0}^{k-1} \frac{\lambda^l}{l!}e^{-\lambda} = \sum_{l=k}^\infty \frac{\lambda^l}{l!}e^{-\lambda}
			\end{align*}
		\section{再生性}
			\begin{shadebox}
				$X_1 \sim \PoissonDist{\lambda_1},\;X_2 \sim \PoissonDist{\lambda_2}\;\Rightarrow\;X_1+X_2 \sim \PoissonDist{\lambda_1 + \lambda_2}$
			\end{shadebox}
			\begin{proof}
				\[ P(X_1+X_2 = k) = \sum_{l=0}^k \frac{\lambda_1^l}{l!}e^{-\lambda_1}\frac{\lambda_2^{k-l}}{(k-l)!}e^{-\lambda_2} = \frac{1}{k!}e^{-(\lambda_1+\lambda_2)} \sum_{l=0}^k \frac{k!}{l!(k-l)!}\lambda_1^l\lambda_2^{k-l} = \frac{1}{k!}(\lambda_1+\lambda_2)^k e^{-(\lambda_1+\lambda_2)} \]
			\end{proof}
	\chapter{Beta分布}
		\section{Beta分布に従う確率変数の生成}
			\subsection{\texorpdfstring{$X\sim\GammaDist{\alpha}{1},Y\sim\GammaDist{\beta}{1}\Rightarrow \frac{X}{X+Y}\sim\Beta{\alpha}{\beta}$}{X～Gamma(α,1), Y～Gamma(β,1) ⇒ X/(X+Y)～Beta(α,β)}}
				\begin{proof}
					\quad\par
					$Z=\frac{X}{X+Y}$とおく。
					\begin{align*}
						\PDF{Z}{z} &= \partDerivLong{\Pr{Z\leq z}}{z}{} = \partDerivLong{\Pr{X\leq\frac{z}{1-z}Y}}{z}{} \\
						&= \partDerivLong{\integrate{y=0}{\infty}{\integrate{x=0}{\frac{z}{1-z}y}{ \PDF{\GammaDist{\alpha}{1}}{x}\PDF{\GammaDist{\beta}{1}}{y} }{}{x}}{}{y}}{z}{} \\
						&= \partDerivLong{\integrate{y=0}{\infty}{\integrate{x=0}{\frac{z}{1-z}y}{\frac{x^{\alpha-1}e^{-x}}{\Gamma(\alpha)}\frac{y^{\beta-1}e^{-y}}{\Gamma(\beta)}}{}{x}}{}{y}}{z}{} \\
						&= \frac{1}{\Gamma(\alpha)\Gamma(\beta)}\integrate{y=0}{\infty}{\partDerivLong{\integrate{x=0}{\frac{z}{1-z}y}{x^{\alpha-1}e^{-x}y^{\beta-1}e^{-y}}{}{x}}{z}{}}{}{y} \\
						&= \frac{1}{\Gamma(\alpha)\Gamma(\beta)}\integrate{y=0}{\infty}{y^{\beta-1}e^{-y}\left(\frac{z}{1-z}y\right)^{\alpha-1}\exp\left(-\frac{z}{1-z}y\right)\left(\partDerivLong{\frac{z}{1-z}y}{z}{}\right)}{}{y} \\
						&= \frac{1}{\Gamma(\alpha)\Gamma(\beta)}\integrate{y=0}{\infty}{y^{\beta-1}\exp\left(\frac{-y}{1-z}y\right)\left(\frac{y}{1-z}\right)^{\alpha-1}z^{\alpha-1}\frac{1}{(1-z)^2}}{}{y} \\
						&= \frac{z^{\alpha-1}(1-z)^{\beta-1}}{\Gamma(\alpha)\Gamma(\beta)}\integrate{0}{\infty}{u^{\alpha+\beta-1}e^{-u}}{}{u} \quad\left(u=\frac{y}{1-z}\textrm{とおいた}\right) \\
						&= \frac{\Gamma(\alpha+\beta)}{\Gamma(\alpha)\Gamma(\beta)}z^{\alpha-1}(1-z)^{\beta-1} = \frac{z^{\alpha-1}(1-z)^{\beta-1}}{B(\alpha,\beta)} = \PDF{\Beta{\alpha}{\beta}}{z}
					\end{align*}
				\end{proof}
	\chapter{正規分布}
		\section{裾確率の評価}
			\begin{shadebox}
				標準正規分布の累積分布関数\texorpdfstring{$\Phi(x)$}{Φ(x)}に対して次が成り立つ。
				\[1-\Phi(\alpha) \leq \frac{1}{2}\exp{\left(-\frac{1}{2}\alpha^2\right) \quad \textrm{for }^\forall \alpha \geq 0}\]
			\end{shadebox}
			\begin{proof}
				標準正規分布の確率密度関数を$\phi(\alpha)$とすると、任意の$\alpha>0$に対して次式が成り立つ。
				\[1-\Phi(\alpha) = \integrate{\alpha}{\infty}{\phi(x)}{}{x} \leq \integrate{\alpha}{\infty}{\frac{x}{\alpha}\phi(x)}{}{x} = \frac{1}{\alpha}\left[-\phi(x)\right]_\alpha^\infty = \frac{1}{\alpha}\phi(\alpha)\]
				よって$\alpha\geq 1$に対しては
				\[1-\Phi(\alpha) \leq \phi(\alpha) = \frac{1}{\sqrt{2\pi}}e^{-\frac{1}{2}\alpha^2} \leq \frac{1}{2}e^{-\frac{1}{2}\alpha^2}\]
				$\alpha<1$に対しては別のやり方で示す。$1-\Phi(\alpha) = 1/2 - \integrate{0}{\alpha}{\phi(x)}{}{x} =: 1/2-S(\alpha)$であり、$\phi(x)$の変曲点は$x=\pm 1$である。
				\begin{figure}[H]
					\centering
					\includegraphics[keepaspectratio, scale=0.5]
					{\currfiledir/imgs/evaluation-of-tail-weight-of-gaussian-distribution/graph.eps}
					\caption{$\phi(x)$のグラフ}
				\end{figure}
				\begin{align*}
					S(\alpha) &\geq \textrm{赤い台形の面積} = \frac{\alpha}{2\sqrt{2\pi}}(1+e^{-\frac{1}{2}\alpha^2}) \\
					\therefore\;1-\Phi(\alpha) &\leq \frac{1}{2}\left(1-\frac{\alpha}{\sqrt{2\pi}}(1+e^{-\frac{1}{2}\alpha^2})\right) \\
					\therefore\;1-\Phi(\alpha)-\frac{1}{2}e^{-\frac{1}{2}\alpha^2} &\leq \frac{1}{2}\left(1-\frac{\alpha}{\sqrt{2\pi}}(1+e^{-\frac{1}{2}\alpha^2})-e^{-\frac{1}{2}\alpha^2}\right) =: g(\alpha)/2
				\end{align*}
				$g(\alpha)$は区間$[0,1]$で凸であり(凸関数の和が凸関数であることと、2階導関数の計算によりわかる)、$g(0)=0$である。
				また、数値計算によると$g(1)\approx 1-\frac{4.8}{4.2}<0$であるから区間$[0,1]$で$g(\alpha)\leq 0$
			\end{proof}
		\section{スケール変換とシフト}
			\begin{shadebox}
				\[X \sim N(\mu,\sigma^2) \Leftrightarrow Z=aX+b \sim N(a\mu+b,(a\sigma)^2) \hspace{10pt} (a \neq 0)\]
			\end{shadebox}
			\begin{proof}
				\quad\par
				$(\Rightarrow)$\newline
				$X$の確率密度関数を$f^X(x)$と書くことにすると
				\[f^X(x) = \frac{1}{\sqrt{2\pi}\sigma}\exp\left[-\frac{(x-\mu)^2}{2\sigma^2}\right]\]
				(i) $a>0$の場合\par
				\[P(-\infty \leq Z \leq z) = P\left(\-\infty \leq X \leq \frac{z-b}{a}\right) = \int_{-\infty}^{\frac{z-b}{a}}{f^X(x) dx}\]
				\[\therefore \; f^Z(z) = \frac{d}{dz}P(-\infty \leq Z \leq z) = f^X(\frac{z-b}{a})\cdot\frac{d}{dz}\left(\frac{z-b}{a}\right) = \frac{1}{\sqrt{2\pi}(a\sigma)}\exp\left[-\frac{(z-a\mu-b)^2}{2(a\sigma)^2}\right]\]
				\[\therefore f^Z(z) = \frac{1}{\sqrt{2\pi}(a\sigma)}\exp\left[-\frac{(z-a\mu-b)^2}{2(a\sigma)^2}\right] = \text{pdf of  }N(a\mu+b,(a\sigma)^2)\]
				(ii) $a<0$の場合\par
				\[P(-\infty \leq Z \leq z) = P\left(\frac{z-b}{a} \leq X \leq \infty\right) = \int_{\frac{z-b}{a}}^{\infty}{f(x) dx}\]
				\begin{align*}
					\therefore \; f^Z(z) &=& \frac{d}{dz}P(-\infty \leq Z \leq z) = \textcolor{red}{-}f^X(\frac{z-b}{a})\cdot\frac{d}{dz}\left(\frac{z-b}{a}\right) = \frac{\textcolor{red}{-1}}{\sqrt{2\pi}(a\sigma)}\exp\left[-\frac{(z-a\mu-b)^2}{2(a\sigma)^2}\right]\\
					&=& \frac{1}{\sqrt{2\pi}(\textcolor{red}{|a|}\sigma)}\exp\left[-\frac{(z-a\mu-b)^2}{2(a\sigma)^2}\right] = \text{pdf of  }N(a\mu+b,(a\sigma)^2)
				\end{align*}
				\newline
				$(\Leftarrow)$\newline
				$Y \sim N(a\mu+b,(b\sigma)^2)$として、$(Y-b)/a\;(=X)$の分布を既に示された必要性に依って計算すれば明らかである。
			\end{proof}
		\section{再生性}
			\begin{shadebox}
				\[X \sim N(\mu_X,{\sigma_X}^2)\;\;,\;\;Y \sim N(\mu_Y,{\sigma_Y}^2) \quad \Rightarrow \quad Z=X+Y \sim N(\mu_X+\mu_Y,{\sigma_X}^2+{\sigma_Y}^2)\]
			\end{shadebox}
			\begin{proof}
				\quad\par
				$X,Y,Z$の確率密度関数を$f^X(x),f^Y(y),f^Z(z)$と書くことにすると
				\begin{align*}
					f^Z(z) &= \derivLong{\integrate{x=-\infty}{\infty}{\integrate{y=-\infty}{z-x}{f^X(x)f^Y(y)}{}{y}}{}{x}}{z}{} = \int_{-\infty}^{\infty}{f^X(x)f^Y(-x+z) dx} \quad(\because\ref{重積分の微分の系}) \\
					&= \int_{-\infty}^{\infty}{\frac{1}{\sqrt{2\pi}\sigma_X}\exp{\left[-\frac{(x-\mu_X)^2}{2{\sigma_X}^2}\right]}\frac{1}{\sqrt{2\pi}\sigma_Y}\exp{\left[-\frac{(-x+z-\mu_Y)^2}{2{\sigma_Y}^2}\right]} dx}\\
					&= \frac{1}{2\pi\sigma_X\sigma_Y}\int_{-\infty}^{\infty}{\exp{\left[-\frac{(x-\mu_X)^2}{2{\sigma_X}^2} - \frac{(-x+z-\mu_Y)^2}{2{\sigma_Y}^2}\right]} dx}
				\end{align*}
				expの引数を$A$とおくと
				\begin{align*}
					A &= -\frac{1}{{2\sigma_X}^2{\sigma_Y}^2}\left[{\sigma_Y}^2(x-\mu_X)^2 + {\sigma_X}^2(-x+z-\mu_Y)^2\right]\\
					&= -\frac{1}{{2\sigma_X}^2{\sigma_Y}^2}\left[{\sigma_X}^2(x^2+{\mu_X}^2-2\mu_Xx)+{\sigma_X}^2(x^2+z^2+{\mu_Y}^2-2zx-2z\mu_Y+2\mu_yx)\right]\\
					&= -\frac{1}{{2\sigma_X}^2{\sigma_Y}^2}\left[({\sigma_X}^2+{\sigma_Y}^2)x^2 + (-2{\sigma_Y}^2\mu_X - 2{\sigma_X}^2z + 2{\sigma_X}^2\mu_Y)x \right.\\
					&\quad \left. + {\sigma_Y}^2{\mu_X}^2 + {\sigma_X}^2z^2 + {\sigma_X}^2{\mu_Y}^2 - 2{\sigma_X}^2z\mu_Y\right]\\
					&= -\frac{1}{{2\sigma_X}^2{\sigma_Y}^2}\left\{{\left[\sqrt{{\sigma_X}^2+{\sigma_Y}^2}x - \frac{{\sigma_Y}^2\mu_X + {\sigma_X}^2z - {\sigma_X}^2\mu_Y}{\sqrt{{\sigma_X}^2+{\sigma_Y}^2}}\right]}^2\right.\\
					&\quad -\frac{1}{{\sigma_X}^2+{\sigma_Y}^2}\left({\sigma_Y}^4{\mu_X}^2 + {\sigma_X}^4z^2 + {\sigma_X}^4{\mu_Y}^2 + 2{\sigma_X}^2{\sigma_Y}^2\mu_X z - 2{\sigma_X}^4\mu_Y z - 2{\sigma_X}^2{\sigma_Y}^2\mu_X\mu_Y\right)\\
					&\quad \left. + {\sigma_Y}^2{\mu_X}^2 + {\sigma_X}^2z^2 + {\sigma_X}^2{\mu_Y}^2 - 2{\sigma_X}^2z\mu_Y \right\}\\
					&= -\frac{1}{{2\sigma_X}^2{\sigma_Y}^2} \left\{{\left[\sqrt{{\sigma_X}^2+{\sigma_Y}^2}x - \frac{{\sigma_Y}^2\mu_X + {\sigma_X}^2z - {\sigma_X}^2\mu_Y}{\sqrt{{\sigma_X}^2+{\sigma_Y}^2}}\right]}^2\right.\\
					&\quad \left. \frac{{\sigma_X}^2{\sigma_Y}^2}{{\sigma_X}^2+{\sigma_Y}^2}({\mu_X}^2 + z^2 + {\mu_Y}^2 - 2\mu_Y z - 2\mu_X z + 2\mu_X\mu_Y) \right\}\\
					&= -\frac{1}{{2\sigma_X}^2{\sigma_Y}^2} \left\{{\left[\sqrt{{\sigma_X}^2+{\sigma_Y}^2}x - \frac{{\sigma_Y}^2\mu_X + {\sigma_X}^2z - {\sigma_X}^2\mu_Y}{\sqrt{{\sigma_X}^2+{\sigma_Y}^2}}\right]}^2 + \frac{{\sigma_X}^2{\sigma_Y}^2}{{\sigma_X}^2+{\sigma_Y}^2}(z-\mu_X-\mu_Y)^2 \right\}
				\end{align*}
				これと(1)より
				\begin{align*}
					f^z(z) &= \frac{1}{2\pi\sigma_X\sigma_Y}\exp{\left[-\frac{(z-\mu_X-\mu_Y)^2}{2({\sigma_X}^2+{\sigma_Y}^2)}\right]}\int_{-\infty}^{\infty}{\exp{\left[\frac{{\sigma_X}^2+{\sigma_Y}^2}{2{\sigma_X}^2{\sigma_Y}^2}\left(x-\frac{{\sigma_Y}^2\mu_X + {\sigma_X}^2z - {\sigma_X}^2\mu_Y}{{\sigma_X}^2+{\sigma_Y}^2}\right)^2\right]} dx}\\
					&= \frac{1}{\sqrt{2\pi}\sqrt{{\sigma_X}^2+{\sigma_Y}^2}}\exp{\left[-\frac{(z-\mu_X-\mu_Y)^2}{2({\sigma_X}^2+{\sigma_Y}^2)}\right]} = \text{pdf of} \quad N(\mu_X+\mu_Y,{\sigma_X}^2+{\sigma_Y}^2)
				\end{align*}
			\end{proof}
		\section{一次結合}
			\begin{shadebox}
				\[X_i \sim N(\mu_i,{\sigma_i}^2) \quad (i=1,2,\dotsc) \quad \Rightarrow Z = \sum_{i=1}^{n}{c_i X_i} \quad \sim N\left(\sum_{i=1}^{n}{c_i \mu_i},\sum_{i=1}^{n}{(c_i \sigma_i)^2}\right)\]
			\end{shadebox}
			\begin{proof}
				\quad\par
				定理1(スケール変換＆シフト),2(和)より明らか。
			\end{proof}
		\section{分散\texorpdfstring{$\sigma^2$}{σ\string^2},平均\texorpdfstring{$\mu$}{μ}のときの累積分布関数は\texorpdfstring{$\Phi(\frac{z-\mu}{\sigma})$}{Φ((z-μ)/σ)}}
			\begin{shadebox}
				標準正規分布の累積分布関数を$\Phi(z) \coloneqq \integrate{-\infty}{z}{\phi(x)}{}{x}$とするとき、
				$X\sim N(\mu,\sigma)$の累積分布関数は$\Phi_X(Z) = \Phi(\frac{z-\mu}{\sigma})$となる。
			\end{shadebox}
			\begin{proof}
				\begin{align*}
					\Phi_X(z) &= \integrate{-\infty}{z}{\frac{1}{\sigma}\func{\phi}{\frac{z-\mu}{\sigma}}}{}{x} = \integrate{-\infty}{\frac{z-\mu}{\sigma}}{}{\phi(y)}{y}\quad\left(\frac{z-\mu}{\sigma}=y\textrm{とおいた}\right) \\
					&= \func{\Phi}{\frac{z-\mu}{\sigma}}
				\end{align*}
			\end{proof}
		\section{零平均,同分散の正規分布に独立に従う\texorpdfstring{$n$}{n}個の確率変数の直交変換もまた同じ分布に従う}
			\label{零平均,同分散の正規分布に独立に従うn個の確率変数の直交変換もまた同じ分布に従う}
			\begin{shadebox}
				$X_1,\dotsc,X_n\sim N(0,\sigma^2),\quad\bm{X}\coloneqq[X_1,\dotsc,X_n]^\top$とし、$P\in\realNumbers^{n\times n}$を直交行列とする。
				$\bm{Y} \coloneqq [Y_1,\dotsc,Y_n]^\top \coloneqq P\bm{X}$とすると、$Y_1,\dotsc,Y_n\sim N(0,\sigma^2)$
			\end{shadebox}
			\begin{proof}
				\quad\par
				$\bm{X}$と$\bm{Y}$の確率密度関数が一致することを示せば良い。
				両者の確率密度関数を各々$p_{\bm{X}},p_{\bm{Y}}$とする。
				任意の点$\bm{y}_0\in\realNumbers^n$における$p_{\bm{Y}}$を計算するために、$\Omega$を、$\bm{y}_0$を含む任意の領域とする。
				$\Omega$を直交変換$P$の逆写像で写した領域を$P^\top(\Omega)$と表記すると
				\begin{align}
					p_{\bm{Y}}(\bm{y}_0) &\coloneqq \lim_{|\Omega|\to 0}\frac{\Pr{\bm{y}\in\Omega}}{|\Omega|} = \lim_{|\Omega|\to 0}\frac{\Pr{\bm{x}\in P^\top(\Omega)}}{|\Omega|} = \lim_{|\Omega|\to 0}\frac{1}{|\Omega|}\integrate{P^\top(\Omega)}{}{p_{\bm{X}}(\bm{x})}{}{\bm{x}} \nonumber\\
					&= \lim_{|\Omega|\to 0}\frac{1}{|\Omega|}\integrate{P^\top(\Omega)}{}{ \frac{1}{(2\pi\sigma^2)^{n/2}}\exp{\left(-\frac{\norm{\bm{x}}^2}{2\sigma^2}\right)} }{}{\bm{x}} \nonumber\\
					&= \lim_{|\Omega|\to 0}\frac{1}{|\Omega|}\integrate{P(P^\top(\Omega))}{}{ \frac{1}{(2\pi\sigma^2)^{n/2}}\exp{\left(-\frac{\norm{P^\top\bm{y}}^2}{2\sigma^2}\right)} |\det{P^\top}|}{}{\bm{y}} \\
					&\quad \quad(\text{↑重積分の変数変換}\bm{x}=P^\top\bm{y}\textrm{を用いた}) \nonumber\\
					&= \lim_{|\Omega|\to 0}\frac{1}{|\Omega|}\integrate{\Omega}{}{ \frac{1}{(2\pi\sigma^2)^{n/2}}\exp{\left(-\frac{\norm{\bm{y}}^2}{2\sigma^2}\right)} }{}{\bm{y}} \quad(P\textrm{は直交行列だから}\det{P^\top}=\pm 1,\;\norm{P^\top\bm{y}}=\norm{\bm{y}}) \nonumber\\
					&= \lim_{|\Omega|\to 0}\frac{1}{|\Omega|}\integrate{\Omega}{}{ p_{\bm{X}}(\bm{y}) }{}{\bm{y}}
				\end{align}
				重積分の平均値の定理より、ある$\tilde{\bm{y}}\in\Omega$が存在して
				\[(1) = \lim_{|\Omega|\to 0}\frac{1}{|\Omega|}|\Omega|p_{\bm{X}}(\tilde{\bm{y}}) = \lim_{|\Omega|\to 0}p_{\bm{X}}(\tilde{\bm{y}}) = p_{\bm{X}}(\bm{y}_0)\]
			\end{proof}
		\section{多変量正規分布}
			\subsection{定義}
				$\Sigma$を正定値対称行列とし、対角化が$\diag{{\sigma_1}^2,\dotsc,{\sigma_n}^2}$であるとする。
				確率変数$X_1,\dotsc,X_n$の同時確率密度関数$f_{\bm{X}}(\bm{x})$が
				\[ f_{\bm{X}}(\bm{x}) = \frac{1}{(2\pi)^{n/2}|\Sigma|^{1/2}} \exp\left[ -\frac{1}{2}(\bm{x}-\bm{\mu})^\top\Sigma^{-1}(\bm{x}-\bm{\mu}) \right]\]
				であるとき、$X_1,\dotsc,X_n$が従う分布は\textbf{多変量正規分布}と呼ばれ、$N(\bm{\mu},\Sigma)$と表記される。
				期待値と共分散行列は$\E{}{\bm{X}}=\bm{\mu},\;\CovMat{\bm{X}} = \Sigma$である。
			\subsection{密度関数の全空間積分が1になることの確認}
				\label{subsec:多変量正規分布の密度関数の全空間積分が1になることの確認}
				$\Sigma$に関する仮定より、適当な直交行列$V$が存在して$\Sigma = V\Lambda V^\top,\;\Lambda=\diag{{\sigma_1}^2,\dotsc,{\sigma_n}^2}$と変形できる。
				そこで$\bm{x}-\bm{\mu} = V\bm{y}$すなわち$\bm{y} = V^\top(\bm{x}-\bm{\mu})$なる直交変換を行うと、求めたい積分は次のように変形できる。
				\begin{align*}
					\integrate{\realNumbers^n}{}{f_{\bm{X}}(\bm{x})}{}{\bm{x}} &= \integrate{\realNumbers^n}{}{\frac{1}{(2\pi)^{n/2}|\Sigma|^{1/2}}\exp\left(-\frac{1}{2}\bm{y}^\top V (V\Lambda V^\top)^{-1} V\bm{y}\right) |\det V|}{}{\bm{y}} \\
					&= \integrate{\realNumbers^n}{}{\frac{1}{\prod_{i=1}^n \sqrt{2\pi}\sigma_i}\exp\left(-\frac{1}{2}\bm{y}^\top\Lambda^{-1}\bm{y}\right)}{}{\bm{y}} \quad (\because\;\text{は直交行列なので}|\det V|=1) \\
					&= \integrate{\realNumbers^n}{}{\prod_{i=1}^n\frac{1}{\sqrt{2\pi}\sigma_i}\exp\left(-\frac{1}{2{\sigma_i}^2}{y_i}^2\right)}{}{\bm{y}} = \prod_{i=1}^n \integrate{-\infty}{\infty}{\frac{1}{\sqrt{2\pi}\sigma_i}\exp\left(-\frac{1}{2{\sigma_i}^2}{y_i}^2\right)}{}{y_i} = 1
				\end{align*}
			\subsection{期待値の導出}
				\begin{proof}
					\quad\par
					先程の直交変換をここでも用いる。
					\begin{align*}
						\E{}{\bm{X}} &= \integrate{\realNumbers^n}{}{\bm{x}f_{\bm{X}}(\bm{x})}{}{\bm{x}} = \integrate{\realNumbers^n}{}{(V\bm{y}+\bm{\mu})\frac{1}{\prod_{i=1}^n \sqrt{2\pi}\sigma_i}\exp\left(-\frac{1}{2}\bm{y}^\top\Lambda^{-1}\bm{y}\right)}{}{\bm{y}} \\
						&= \bm{\mu} + V\integrate{\realNumbers^n}{}{\bm{y}\frac{1}{\prod_{i=1}^n \sqrt{2\pi}\sigma_i}\exp\left(-\frac{1}{2}\bm{y}^\top\Lambda^{-1}\bm{y}\right)}{}{\bm{y}} \\
						&= \bm{\mu} + V\integrate{\realNumbers^n}{}{\bm{y}\prod_{i=1}^n\frac{1}{\sqrt{2\pi}\sigma_i}\exp\left(-\frac{1}{2{\sigma_i}^2}{y_i}^2\right)}{}{\bm{y}} \\
						&= \bm{\mu} + V\integrate{y_n=-\infty}{\infty}{\dotsi\integrate{y_1=-\infty}{\infty}{\bm{y}\prod_{i=1}^n\frac{1}{\sqrt{2\pi}\sigma_i}\exp\left(-\frac{1}{2{\sigma_i}^2}{y_i}^2\right)}{}{y_1}\dotsi}{}{y_n}
					\end{align*}
					第二項目はベクトルの各インデックス毎に奇関数の対称積分が起こり、結局全体が0になることがわかる。
				\end{proof}
			\subsection{共分散行列の導出}
				\begin{proof}
					\quad\par
					\ref{subsec:多変量正規分布の密度関数の全空間積分が1になることの確認}で用いた直交変換をここでも用いる。
					\begin{align*}
						\CovMat{\bm{X}} &= \E{}{(\bm{X}-\E{}{\bm{X}})(\bm{X}-\E{}{\bm{X}})^\top} = \integrate{\realNumbers^n}{}{(\bm{x}-\bm{\mu})(\bm{x}-\bm{\mu})^\top f_{\bm{X}}(\bm{x})}{}{\bm{x}} \\
						&= \integrate{\realNumbers^n}{}{V\bm{y}\bm{y}^\top V^\top\prod_{i=1}^n\frac{1}{\sqrt{2\pi}\sigma_i}\exp\left(-\frac{1}{2{\sigma_i}^2}{y_i}^2\right)}{}{\bm{y}} \\
						&= V\left(\integrate{\realNumbers^n}{}{\bm{y}\bm{y}^\top \prod_{i=1}^n\frac{1}{\sqrt{2\pi}\sigma_i}\exp\left(-\frac{1}{2{\sigma_i}^2}{y_i}^2\right)}{}{\bm{y}}\right)V^\top
					\end{align*}
					積分の部分に着目すると、非対角要素では奇関数の対称積分が起こり、0になる。
					第$k$対角成分は、$y_i(i\neq k)$に関する積分を先に実行することで
					\begin{align*}
						\integrate{\realNumbers^n}{}{{y_k}^2\prod_{i=1}^n\frac{1}{\sqrt{2\pi}\sigma_i}\exp\left(-\frac{1}{2{\sigma_i}^2}{y_i}^2\right)}{}{\bm{y}} = \integrate{-\infty}{\infty}{{y_k}^2\frac{1}{\sqrt{2\pi}\sigma_k}\exp\left(-\frac{1}{2{\sigma_k}^2}{y_k}^2\right)}{}{y_k} = {\sigma_k}^2
					\end{align*}
					となることがわかる。
					よって
					\[ \CovMat{\bm{X}} = V\Lambda V^\top = \Sigma \]
				\end{proof}
			\subsection{特性関数は\texorpdfstring{$\phi_{\bm{X}}(\bm{t})=\exp(i\bm{t}^\top\bm{\mu} - \bm{t}^\top\Sigma\bm{t}/2)$}{φ\_X (t) = exp(i t\string^T μ - t\string^T Σ t/2)}}
				\begin{proof}
					\quad\par
					\ref{subsec:多変量正規分布の密度関数の全空間積分が1になることの確認}で用いた直交変換をここでも用いる。
					\begin{align*}
						\phi_{\bm{X}}(\bm{t}) &= \E{}{\exp(i\bm{t}^\top \bm{X})} = \integrate{\realNumbers^n}{}{\exp(i\bm{t}^\top\bm{x})f_{\bm{X}}(\bm{x})}{}{\bm{x}} \\
						&= \integrate{\realNumbers^n}{}{\frac{1}{(2\pi)^{n/2}|\Sigma|^{1/2}}\exp(i\bm{t}^\top\bm{x})\exp\left(-\frac{1}{2}(\bm{x}-\bm{\mu})^\top\Sigma^{-1}(\bm{x}-\bm{\mu})\right)}{}{\bm{x}} \\
						&= \integrate{\realNumbers^n}{}{\frac{1}{(2\pi)^{n/2}\prod_{i=1}^n\sigma_i}\exp(i\bm{t}^\top(V\bm{y}+\bm{\mu}))\exp\left(-\frac{1}{2}\bm{y}^\top\Lambda^{-1}\bm{y}\right)|\det V|}{}{\bm{y}} \\
						&= \exp(i\bm{t}^\top\bm{\mu})\integrate{\realNumbers^n}{}{\frac{1}{(2\pi)^{n/2}\prod_{i=1}^n\sigma_i}\exp\left(-\frac{1}{2}\bm{y}^\top\Lambda^{-1}\bm{y} + i\bm{t}^\top V\bm{y}\right)}{}{\bm{y}} \\
						&= \exp(i\bm{t}^\top\bm{\mu})\integrate{\realNumbers^n}{}{\frac{1}{(2\pi)^{n/2}\prod_{i=1}^n\sigma_i}\exp\left(-\frac{1}{2}\sum_{i=1}^n\frac{1}{{\sigma_i}^2}{y_i}^2 + i\bm{t}^\top V\bm{y}\right)}{}{\bm{y}}
					\end{align*}
					ここで$c_i \coloneqq (\bm{t}^\top V)[i]$とするとexpの引数は次のように変形できる。
					\begin{align*}
						& -\frac{1}{2}\sum_{i=1}^n\frac{1}{{\sigma_i}^2}{y_i}^2 + i\bm{t}^\top V\bm{y} = -\frac{1}{2}\sum_{i=1}^n\frac{1}{{\sigma_i}^2}{y_i}^2 + i\sum_{i=1}^n c_iy_i = -\frac{1}{2}\sum_{i=1}^n\frac{1}{{\sigma_i}^2}\left({y_i}^2 - i2{\sigma_i}^2c_iy_i\right) \\
						=& -\frac{1}{2}\sum_{i=1}^n\frac{1}{{\sigma_i}^2}\left[\left(y_i - i{\sigma_i}^2c_i\right)^2 + {\sigma_i}^4{c_i}^2\right] = -\frac{1}{2}\sum_{i=1}^n{\sigma_i}^2{c_i}^2 - \frac{1}{2}\sum_{i=1}^n\frac{1}{{\sigma_i}^2}\left(y_i - i{\sigma_i}^2c_i\right)^2
					\end{align*}
					さらに次式が成り立つ。
					\[ \sum_{i=1}^n{\sigma_i}^2{c_i}^2 = [c_1,\dotsc,c_n]\Lambda[c_1,\dotsc,c_n]^\top = \bm{t}^\top V\Lambda V^\top\bm{t} = \bm{t}^\top\Sigma\bm{t} \]
					よって
					\begin{align*}
						\phi_{\bm{X}}(t) &= \exp\left(i\bm{t}^\top\bm{\mu} - \frac{1}{2}\bm{t}^\top\Sigma\bm{t}\right)\integrate{\realNumbers^n}{}{\frac{1}{(2\pi)^{n/2}\prod_{i=1}^n\sigma_i}\exp\left(-\frac{1}{2}\sum_{i=1}^n\frac{1}{{\sigma_i}^2}\left(y_i - i{\sigma_i}^2c_i\right)^2\right)}{}{\bm{y}} \\
						&= \exp\left(i\bm{t}^\top\bm{\mu} - \frac{1}{2}\bm{t}^\top\Sigma\bm{t}\right)\integrate{\realNumbers^n}{}{\prod_{i=1}^n\frac{1}{\sqrt{2\pi}\sigma_i}\exp\left(-\frac{1}{2{\sigma_i}^2}\left(y_i - i{\sigma_i}^2c_i\right)^2\right)}{}{\bm{y}} \\
						&= \exp\left(i\bm{t}^\top\bm{\mu} - \frac{1}{2}\bm{t}^\top\Sigma\bm{t}\right)\prod_{i=1}^n I_i\quad \left(I_i \coloneqq \integrate{-\infty}{\infty}{\frac{1}{\sqrt{2\pi}\sigma_i}\exp\left(-\frac{1}{2{\sigma_i}^2}\left(y_i - i{\sigma_i}^2c_i\right)^2\right)}{}{y_i}\right)
					\end{align*}
					$I_i$は複素積分により1である(矩形閉路ABCDA: $A=(-L,0),\;B=(L,0),\;C=(L,i{\sigma_i}^2c_i),\;D=(-L,i{\sigma_i}^2c_i)$を用いて$L\to\infty$とするとよい)ので
					\[ \phi_{\bm{X}}(t) = \exp\left(i\bm{t}^\top\bm{\mu} - \frac{1}{2}\bm{t}^\top\Sigma\bm{t}\right) \]
				\end{proof}
			\subsection{特性関数を用いた期待値と分散の導出}
				\begin{proof}
					\quad\par
					まず、期待値については$\nabla_{\bm{t}}\phi_{\bm{X}}(\bm{t}) = (i\bm{\mu}-\Sigma\bm{t})\phi_{\bm{X}}(\bm{t})$より$\E{}{\bm{X}} = -i\nabla_{\bm{t}}\phi_{\bm{X}}(\bm{0}) = \bm{\mu}$である。
					次に、共分散行列については、
					\[ \nabla_{\bm{t}}^2\phi_{\bm{X}}(\bm{t}) = (-\Sigma + (i\bm{\mu}-\Sigma\bm{t})(i\bm{\mu}-\Sigma\bm{t})^\top)\phi_{\bm{X}}(\bm{t}) \]
					より$\E{}{\bm{X}\bm{X}^\top} = -\nabla_{\bm{t}}^2\phi_{\bm{X}}(\bm{0}) = \Sigma + \bm{\mu}\bm{\mu}^\top$であるので
					$\CovMat{\bm{X}} = \E{}{\bm{X}\bm{X}^\top} - \E{}{\bm{X}}\E{}{\bm{X}}^\top = \Sigma$
				\end{proof}
			\subsection{線形写像 : \texorpdfstring{$A\bm{X}+\bm{b} \sim N(A\bm{\mu}+\bm{b}, A\Sigma A^\top)$}{AX+b ～ N(Aμ+b, AΣA\string^T)}}
				\begin{proof}
					\quad\par
					特性関数を用いて証明する。
					\begin{align*}
						\varphi_{A\bm{X}+\bm{b}}(\bm{t}) &= \E{}{\exp(i\bm{t}^\top(A\bm{X}+\bm{b}))} = \exp(i\bm{t}^\top\bm{b})\E{}{\exp(i\bm{t}^\top A\bm{X})} = \exp(i\bm{t}^\top\bm{b})\E{}{\exp(i(A^\top\bm{t})^\top\bm{X})} \\
						&= \exp(i\bm{t}^\top\bm{b})\varphi_{\bm{X}}(A^\top\bm{t}) = \exp(i\bm{t}^\top\bm{b})\exp\left(i(A^\top\bm{t})^\top\bm{\mu} - \frac{1}{2}(A^\top\bm{t})^\top\Sigma(A^\top\bm{t})\right) \\
						&= \exp(i\bm{t}^\top\bm{b})\exp\left(i\bm{t}^\top(A\bm{\mu}) - \frac{1}{2}\bm{t}^\top(A\Sigma A^\top)\bm{t}\right) = \exp\left(i\bm{t}^\top(A\bm{\mu}+\bm{b}) - \frac{1}{2}\bm{t}^\top(A\Sigma A^\top)\bm{t}\right)
					\end{align*}
					これは期待値$A\bm{\mu}+\bm{b}$,\;共分散行列$A\Sigma A^\top$の多変量正規分布の特性関数であるから、$A\bm{X}+\bm{b}$がこの多変量正規分布に従うことがわかる。
				\end{proof}
			\subsection{周辺分布 : \texorpdfstring{$X_i \sim N(\mu_i, {\sigma_i}^2)$}{Xi～N(μi, (σi)\string^2)}}
				\begin{proof}
					\quad 直前の定理で$A=\bm{e}_i^\top,\;\bm{b}=0$とすればよい。
				\end{proof}
	\chapter{対数正規分布}
		\section{確率密度関数の導出}
			\begin{shadebox}
				対数正規分布$\Lambda(\mu,\sigma^2)$の確率密度関数$f$は次式である。
				\[ f(x) = \frac{1}{\sqrt{2\pi}\sigma x_2}\exp\left(-\frac{(\log x-\mu)^2}{2\sigma^2}\right) \]
			\end{shadebox}
			確率変数$X>0$の対数$Y=\log X$が期待値$\mu$, 分散$\sigma$の正規分布$N(\mu,\sigma^2)$に従うという条件から$X$の確率密度関数$f_X$を導出する。
			$x_1,x_2>0$とすると次式が成り立つ。
			\[ \Pr{x_1\leq X<x_2} = \Pr{\log x_1 \leq Y<\log x_2} = \integrate{\log x_1}{\log x_2}{f_Y(y)}{}{y} \]
			ここに$f_Y$は$N(\mu,\sigma^2)$の確率密度関数である。
			$f_X$は次式となる。
			\begin{align*}
				f_X(x_2) &= \partDerivLong{\Pr{x_1\leq X<x_2}}{x_2}{} = f_Y(\log x_2)\derivLong{\log x_2}{x_2}{} = \frac{1}{x_2}f_Y(\log x_2) \\
				&= \frac{1}{\sqrt{2\pi}\sigma x_2}\exp\left(-\frac{(\log x_2-\mu)^2}{2\sigma^2}\right)
			\end{align*}
		\section{最頻値}
			\begin{shadebox}
				対数正規分布$\Lambda(\mu,\sigma^2)$の最頻値は$e^{\mu-\sigma^2}$である。
			\end{shadebox}
			\begin{proof}
				\quad\par
				対数正規分布の確率密度関数$f$が1峰性であることは既知として、最頻値すなわち$f$の導関数が0となる点を求める。
				\begin{align*}
					\deriv{f_X(x)}{x}{} = -\frac{1}{\sqrt{2\pi}\sigma x^2}\exp\left(-\frac{(\log x_2-\mu)^2}{2\sigma^2}\right)\left(1+\frac{\log x-\mu}{\sigma^2}\right)
				\end{align*}
				これが0となる$x$の値は$e^{\mu-\sigma^2}$である。
			\end{proof}
		\section{期待値}
			\begin{shadebox}
				対数正規分布$\Lambda(\mu,\sigma^2)$の期待値は$e^{\mu+\sigma^2/2}$である。
			\end{shadebox}
			\begin{proof}
				\quad\par
				$X\sim\Lambda(\mu,\sigma^2)$とする。
				$X$の確率密度関数を$f$とする。
				\begin{align*}
					\E{}{X} &= \integrate{+0}{\infty}{xf(x)}{}{x} = \integrate{+0}{\infty}{\frac{1}{\sqrt{2\pi}\sigma}\exp\left(-\frac{(\log x-\mu)^2}{2\sigma^2}\right)}{}{x} \\
					&= \integrate{-\infty}{\infty}{\frac{1}{\sqrt{2\pi}\sigma}\exp\left(-\frac{(y-\mu)^2}{2\sigma^2}\right)e^y}{}{y} \quad (\text{変数変換 }y=\log x) \\
					&=  \integrate{-\infty}{\infty}{\frac{1}{\sqrt{2\pi}\sigma}\exp\left(-\frac{(y-\mu-\sigma^2)^2}{2\sigma^2} + \mu+\sigma^2/2\right)}{}{y} \\
					&= e^{\mu+\sigma^2/2}\integrate{-\infty}{\infty}{\frac{1}{\sqrt{2\pi}\sigma}\exp\left(-\frac{(y-\mu-\sigma^2)^2}{2\sigma^2}\right)}{}{y} = e^{\mu+\sigma^2/2}
				\end{align*}
			\end{proof}
	\chapter{\texorpdfstring{$\chi^2$}{χ\string^2}分布}
		\section{定義}
			\begin{shadebox}
				関数$f(x)$を次のように定義する。
				\[f(x) \coloneqq \frac{1}{2^{\frac{n}{2}\Gamma\left(\frac{n}{2}\right)}}e^{-\frac{x}{2}}x^{\frac{n}{2}-1}\]
				確率密度関数が上式で与えられるような分布を「自由度$n$の$\chi^2$分布」と呼び、${\chi^2}_n$と書くことにする。
			\end{shadebox}
		\section{\texorpdfstring{$X \sim N(0,1) \quad \Rightarrow \quad X^2 \sim \chi^2_1$}{X～N(0,1) ⇒ X\string^2～χ\string^2\_1}}
			\begin{shadebox}
				\[X \sim N(0,1) \quad \Rightarrow \quad X^2 \sim \chi^2_1\]
			\end{shadebox}
			\begin{proof}
				\quad\par
				$X$の確率密度関数を$f^X(x)$と書くことにすると
				\[P(0\leq Z \leq z) = P(-\sqrt{z}\leq X \leq\sqrt{z}) = rate{-\sqrt{z}}{\sqrt{z}}{f^X(x)}{}{x} = 2\integrate{0}{\sqrt{z}}{f^X(x)}{}{x} = \]
				\begin{align*}
					\therefore \; f^Z(z) &= \derivLong{P(0\leq Z \leq z)}{z}{} = \left(\derivLong{P(0\leq Z \leq z)}{\sqrt{z}}{}\right)\deriv{\sqrt{z}}{z}{} = \frac{1}{\sqrt{z}}f^X(\sqrt{z}) = \frac{1}{\sqrt{z}\sqrt{2\pi}}e^{-\frac{z}{2}}\\
					&= \frac{1}{2^{\frac{1}{2}}\Gamma(\frac{1}{2})}e^{-\frac{z}{2}}z^{\frac{1}{2}-1} = \text{pdf of}\;\chi^2_1
				\end{align*}
			\end{proof}
		\section{\texorpdfstring{$X_i \sim N(0,1) \quad \Rightarrow \quad Z_n = \sum_{i=1}^{n}{{X_i}^2} \sim {\chi^2}_n$}{Xi～N(0,1) ⇒ Zn = (X1)\string^2 + ... + (Xn)\string^2 ～ χ\string^2\_n}}
			\label{標準正規分布に従うn個の独立な確率変数の2乗和は自由度n-1のカイ2乗分布に従う}
			\begin{shadebox}
				\[X_i \sim N(0,1) \quad \Rightarrow \quad Z_n = \sum_{i=1}^{n}{{X_i}^2} \sim {\chi^2}_n\]
			\end{shadebox}
			\begin{proof}
				\quad\par
				定理1より$n=1$のとき成り立つ。$n-1$まで成り立つと仮定して$n$のとき成り立つことを示せば良い。\\
				$Z_n$の確率密度関数を$f^{Z_n}(z)$と書くことにする。$n-1$まで成り立つと仮定すると当然
				\[f^{Z_{n-1}}(z) = \frac{1}{2^{\frac{n-1}{2}}\Gamma\left(\frac{n-1}{2}\right)}e^{-\frac{z}{2}}z^{\frac{n-1}{2}-1}\]
				であるから
				\begin{align*}
					f^{Z_n}(z) &= \int_{x=0}^{\infty}{\int_{y=0}^{\infty}{\delta(x+y-z)f^{Z_{n-1}}(x)f^{Z_1}(y) dy} dx} = \int_{x=0}^{\infty}{f^{Z_{n-1}}(x)\left[\int_{y=0}^{\infty}{\delta(x+y-z)f^{Z_1}(y) dy}\right] dx}\\
					&= \int_{x=0}^{\textcolor{red}{z}}{f^{Z_{n-1}}(x)f^{Z_1}(z-x) dx} = \int_{0}^{z}{\frac{e^{-\frac{x}{2}}}{2^{\frac{n-1}{2}}\Gamma\left(\frac{n-1}{2}\right)}x^{\frac{n-1}{2}-1}\frac{e^{\frac{x-z}{2}}}{2^{\frac{1}{2}}\Gamma\left(\frac{1}{2}\right)}(z-x)^{-\frac{1}{2}} dx}\\
					&= \frac{e^{-\frac{z}{2}}}{2^{\frac{n}{2}}\Gamma\left(\frac{n-1}{2}\right)\Gamma\left(\frac{1}{2}\right)}\int_{0}^{z}{x^{\frac{n-1}{2}-1}(z-x)^{\frac{n}{2}-1} dx}
				\end{align*}
				$x=uz$なる変数変換を行って
				\begin{align*}
					&\frac{e^{-\frac{z}{2}}}{2^{\frac{n}{2}}\Gamma\left(\frac{n-1}{2}\right)\Gamma\left(\frac{1}{2}\right)}\int_{0}^{1}{(uz)^{\frac{n-1}{2}-1}z^{-\frac{1}{2}}(1-u)^{-\frac{1}{2}} zdu}\\
					&= \frac{e^{-\frac{z}{2}}}{2^{\frac{n}{2}}\Gamma\left(\frac{n-1}{2}\right)\Gamma\left(\frac{1}{2}\right)}z^{\frac{n}{2}-1}\int_{0}^{1}{u^{\frac{n-1}{2}-1}(1-u)^{\frac{1}{2}-1} du} = \frac{e^{\frac{z}{2}}}{2^{\frac{n}{2}}\Gamma\left(\frac{n-1}{2}\right)\Gamma\left(\frac{1}{2}\right)}z^{\frac{n}{2}-1}B\left(\frac{n-1}{2},\frac{1}{2}\right)\\
					&= \frac{e^{-\frac{z}{2}}}{2^{\frac{n}{2}}\Gamma\left(\frac{n-1}{2}\right)\Gamma\left(\frac{1}{2}\right)}z^{\frac{n}{2}-1}\frac{\Gamma\left(\frac{n-1}{2}\right)\Gamma\left(\frac{1}{2}\right)}{\Gamma\left(\frac{n}{2}\right)} \hspace{10pt} (\because \; (定理1.3.1 : ベータ関数とガンマ関数の関係))\\
					&= \frac{1}{2^{\frac{n}{2}}\Gamma\left(\frac{n}{2}\right)}e^{-\frac{z}{2}}z^{\frac{n}{2}-1} = \text{pdf of} \quad {\chi^2}_n
				\end{align*}
			\end{proof}
		\section{\texorpdfstring{$X \sim {\chi^2}_n \quad \Rightarrow \quad E[X] = n,\;\;\;V[x] = 2n$}{X～χ\string^2\_n ⇒ E(X) = n, V(X) = 2n}}
			\begin{shadebox}
				\[X \sim {\chi^2}_n \quad \Rightarrow \quad E[X] = n,\;\;\;V[x] = 2n\]
			\end{shadebox}
			\begin{proof}
				\quad\par
				$X$の確率密度関数を$f(x)$と書くことにする。\\
				まず
				\[E[X] = \int_{0}^{\infty}{xf(x) dx} = \frac{1}{2^{\frac{n}{2}}\Gamma{\left(\frac{n}{2}\right)}}\int_0^{\infty}{x^{\frac{n}{2}}e^{-\frac{x}{2}} dx}\]
				$\frac{x}{2} = u$なる変数変換を行って
				\begin{align*}
					& \frac{1}{2^{\frac{n}{2}}\Gamma\left(\frac{n}{2}\right)}\int_0^{\infty}{(2u)^{\frac{n}{2}}e^{-u} 2du} = \frac{2}{\Gamma{\left(\frac{n}{2}\right)}}\int_0^{\infty}{u^{\frac{n}{2}}e^{-u} du} = \frac{2}{\Gamma{\left(\frac{n}{2}\right)}}\int_0^{\infty}{u^{\left(\frac{n}{2}+1\right)-1}e^{-u} du}\\
					&= \frac{2\Gamma{\left(\frac{n}{2}+1\right)}}{\Gamma{\left(\frac{n}{2}\right)}} =  \frac{2\frac{n}{2}\Gamma{\left(\frac{n}{2}\right)}}{\Gamma{\left(\frac{n}{2}\right)}} = n
				\end{align*}
				期待値が$n$になることが示された。分散を求めるために2乗の期待値も計算しておく。
				\[E[X^2] = \int_{0}^{\infty}{x^2f(x) dx} = \frac{1}{2^{\frac{n}{2}}\Gamma{\left(\frac{n}{2}\right)}}\int_0^{\infty}{x^{\frac{n}{2}+1}e^{-\frac{x}{2}} dx}\]
				$\frac{x}{2} = u$なる変数変換を行って
				\begin{align*}
					&\frac{1}{2^{\frac{n}{2}}\Gamma{\left(\frac{n}{2}\right)}}\int_0^{\infty}{(2u)^{\frac{n}{2}+1}e^{-u} 2du} = \frac{4}{\Gamma{\left(\frac{n}{2}\right)}}\int_0^{\infty}{u^{\left(\frac{n}{2}+2\right)-1}e^{-u} du}\\
					&= \frac{4}{\Gamma{\left(\frac{n}{2}\right)}}\Gamma{\left(\frac{n}{2}+2\right)} = \frac{4}{\Gamma{\left(\frac{n}{2}\right)}}\left(\frac{n}{2}+1\right)\left(\frac{n}{2}\right)\Gamma{\left(\frac{n}{2}\right)} = n^2+2n
				\end{align*}
				最後に分散は
				\[V[X] = E[X^2] - {E[X]}^2 = n^2+2n-n^2 = 2n\]
			\end{proof}
		\section{再生性}
			\begin{shadebox}
				$X\sim{\chi^2}_{n_X}, Y\sim{\chi^2}_{n_Y} \Rightarrow X+Y\sim{\chi^2}_{n_X+n_Y}$
			\end{shadebox}
			\begin{proof}
				\quad\par
				$Z=X+Y$とする。$X,Y,Z$の確率密度関数を$f^X(x),f^Y(y),f^Z(z)$とすると
				\begin{align*}
					f^Z(z) &= \derivLong{\integrate{x=0}{z}{\integrate{y=0}{z-x}{f^X(x)f^Y(y)}{}{y}}{}{x}}{z}{} = \int_{0}^{z}{f^X(x)f^Y(z-x) dx} \quad(\because\ref{重積分の微分})\\
					&= \int_{0}^{z}{\frac{1}{2^{\frac{n_X}{2}}\Gamma{\left(\frac{n_X}{2}\right)}}x^{\frac{n_X}{2}-1}e^{-\frac{x}{2}}\frac{1}{2^{\frac{n_Y}{2}}\Gamma{\left(\frac{n_Y}{2}\right)}}(z-x)^{\frac{n_Y}{2}-1}e^{-\frac{z-x}{2}} dx}\\
					&= \frac{1}{2^{\frac{n_X+n_Y}{2}}\Gamma{\left(\frac{n_X}{2}\right)}\Gamma{\left(\frac{n_Y}{2}\right)}}e^{-\frac{z}{2}}\int_{0}^{z}{x^{\frac{n_X}{2}-1}(z-x)^{\frac{n_Y}{2}-1} dx}
				\end{align*}
				積分値の係数が煩わしいのでこれを$A$と書くことにし、$u=\frac{x}{z}$なる変数変換を行って
				\begin{align*}
					f^Z(z) &= Az^{\frac{n_X+n_Y}{2}-1}\int_{0}^{1}{u^{\frac{n_X}{2}-1}(1-u)^{\frac{n_Y}{2}-1} du}\\
					&= Az^{\frac{n_X+n_Y}{2}-1}B\left(\frac{n_X}{2},\frac{n_Y}{2}\right)\\
					&= Az^{\frac{n_X+n_Y}{2}-1}\frac{\Gamma{\left(\frac{n_X}{2}\right)}\Gamma{\left(\frac{n_Y}{2}\right)}}{\Gamma{\left(\frac{n_X+n_Y}{2}\right)}}\\
					&= \frac{1}{2^{\frac{n_X+n_Y}{2}}\Gamma{\left(\frac{n_X}{2}\right)}\Gamma{\left(\frac{n_Y}{2}\right)}}e^{-\frac{z}{2}}z^{\frac{n_X+n_Y}{2}-1}\frac{\Gamma{\left(\frac{n_X}{2}\right)}\Gamma{\left(\frac{n_Y}{2}\right)}}{\Gamma{\left(\frac{n_X+n_Y}{2}\right)}}\\
					&= \frac{1}{2^{\frac{n_X+n_Y}{2}}\Gamma{\left(\frac{n_X+n_Y}{2}\right)}}z^{\frac{n_X+n_Y}{2}-1}e^{-\frac{z}{2}}\\
					&= \text{pdf of }{\chi^2}_{n_X+n_Y}
				\end{align*}
			\end{proof}
		\section{\texorpdfstring{$X_1,\dotsc,X_n\sim N(\mu,\sigma^2) \Rightarrow \frac{1}{\sigma^2}\sum_{i=1}^n{(X_i-\overline{X})^2} \sim \chi^2_{n-1}$}{X1, ..., Xn ～ N(μ,σ\string^2) ⇒ (1/σ\string^2)[(X1 - sampleAvg(X))\string^2 + ... + (Xn-E(X))\string^2] ～ χ\string^2\_(n-1)}}
			\begin{shadebox}
				$X_1,\dotsc,X_n\sim N(\mu,\sigma^2)$とし、$\overline{X}$を標本平均とすると$\frac{1}{\sigma^2}\sum_{i=1}^n{(X_i-\overline{X})^2} \sim \chi^2_{n-1}$
			\end{shadebox}
			\begin{proof}
				\quad\par
				$Z_i \coloneqq \frac{X_i-\mu}{\sigma}$とすると$Z_i \sim N(0,1)$であり
				\begin{equation}
					\frac{1}{\sigma^2}\sum_{i=1}^n{(X_i-\overline{X})^2} = \sum_{i=1}^n{(Z_i-\overline{Z})^2} \quad \textrm{where }\overline{Z}\coloneqq\frac{1}{n}\sum_{i=1}^n{X_i}
				\end{equation}
				$\bm{Z} \coloneqq [Z_1,\dotsc,Z_n]^\top$とする。
				最下行が全て$1/\sqrt{n}$である$n$次直交行列を$P$として$\bm{Y} \coloneqq [Y_1,\dotsc,Y_n]^\top \coloneqq P\bm{Z}$とすると、\ref{零平均,同分散の正規分布に独立に従うn個の確率変数の直交変換もまた同じ分布に従う}より、$Y_1,\dotsc,Y_n\sim N(0,1)$であり、$Y_n = \frac{1}{\sqrt{n}}\sum_{i=1}^n{Z_i} = \sqrt{n}\overline{Z} \therefore\;\overline{Z}=Y_n/\sqrt{n}$。
				\[(1) = \sum_{i=1}^n{({Z_i}^2 + \overline{Z}^2 -2\overline{Z}Z_i)} = \norm{\bm{Z}}^2 - n\overline{Z}^2 = \norm{\bm{Y}}^2 - \overline{Y_n}^2 = \sum_{i=1}^{n\textcolor{red}{-1}}{{Y_i}^2}\]
				右から2つ目の等号で、直交変換がノルムを保存する性質を用いた。
				このように、標準正規分布に従う独立な$n-1$個の確率変数の2乗和で表されたので、\ref{標準正規分布に従うn個の独立な確率変数の2乗和は自由度n-1のカイ2乗分布に従う}よりこれは自由度$n-1$のカイ2乗分布に従う。
			\end{proof}
			\begin{commentary}
				\quad\par
				なぜ独立変数の数を減らせたのか考えてみよう。
				元々の$\bm{Z}$では$Z_1,\dotsc,Z_n$が一斉に同じ数だけ増えても$\sum_{i=1}^n{(Z_i-\overline{Z})^2}$の値は変わらない。
				これは、確率密度関数を求めるための重積分中の$\bm{Z}$の拘束領域が、ベクトル$\bm{1}^n$方向に伸び放題であるということだ。
				$\bm{Z}$はこの方向に無制限に動ける。
				上述の直交変換で$Y_n$がこの方向に割り当てられたため、確率密度関数を求めるための重積分の中のこの方向に関する$(-\infty,\infty)$の積分結果が$Y_n$の確率密度関数の全域積分すなわち1になり、我々が相手にするべき次元が1個減ったのである。
			\end{commentary}
	\chapter{\texorpdfstring{$t$}{t}分布}
		\section{定義}
			\begin{shadebox}
				関数$f(t)$を次のように定義する。
				\[f(t) \coloneqq \frac{\Gamma{\left(\frac{n+1}{2}\right)}}{\sqrt{n\pi}\Gamma{\left(\frac{n}{2}\right)}}\left(1+\frac{t^2}{n}\right)^{-\frac{n+1}{2}}\]
				確率密度関数が上式で与えられるような分布を「自由度$n$の$t$分布」と呼び、$t_n$と書くことにする。
			\end{shadebox}
		\section{\thesubsection \texorpdfstring{$X,Y : \text{indep},\;\;\;X \sim N(0,1),\;\;\;Y \sim {\chi^2}_n \quad \Rightarrow \quad T \coloneqq \frac{X}{\sqrt{\frac{Y}{2}}} \sim t_n$}{X,Y: 独立, X～N(0,1), Y～χ\string^2\_n ⇒ T := X/√(Y/2) ～ t\_n}}
			\begin{shadebox}
				\[X,Y : \text{indep},\;\;\;X \sim N(0,1),\;\;\;Y \sim {\chi^2}_n \quad \Rightarrow \quad T \coloneqq \frac{X}{\sqrt{\frac{Y}{2}}} \sim t_n\]
			\end{shadebox}
			\begin{proof}
				\quad\par
				$X,Y,T$の確率密度関数を$f^X(x),f^Y(y),f~T(t)$と書くことにする。
				\begin{align*}
					&\quad P(-\infty \leq T \leq t) = P\left(-\infty \leq X \leq t\sqrt{\frac{y}{n}}\right) = \int_{0}^{\infty}{\int_{-\infty}^{t\sqrt{\frac{y}{n}}}{f^Y(y)f^X(x) dx} dy} \\
					&= \int_{0}^{\infty}{f^Y(y)\left(\int_{-\infty}^{t\sqrt{\frac{y}{n}}}{f^X(x) dx}\right) dy}
				\end{align*}
				\begin{align*}
					\therefore \; f^T(t) &= \frac{\partial}{\partial t}P(-\infty \leq T \leq t) = \int_{0}^{\infty}{f^Y(y)f^X\left(t\sqrt{\frac{y}{n}}\right)\sqrt{\frac{y}{n}} dy}\quad\text{(微分と積分の順序を交換した)}\\
					&=\frac{1}{\sqrt{2n\pi}2^{\frac{n}{2}}\Gamma{\left(\frac{n}{2}\right)}}\int_{0}^{\infty}{y^{\frac{n}{2}-1}e^{-\frac{y}{2}}\exp{\left(-\frac{t^2y}{2n}\right)}y^{\frac{1}{2}} dy} \\
					&= \frac{1}{\sqrt{2n\pi}2^{\frac{n}{2}}\Gamma{\left(\frac{n}{2}\right)}}\int_{0}^{\infty}{y^{\frac{n-1}{2}}\exp{\left[-\frac{y}{2}\left(1+\frac{t^2}{n}\right)\right]} dy}\\
					&\left(\frac{y}{2}\left(1+\frac{t^2}{n}\right) = u \text{なる変数変換を行って}\right)\\
					&=\frac{1}{\sqrt{2n\pi}2^{\frac{n}{2}}\Gamma{\left(\frac{n}{2}\right)}}\left(\frac{2}{1+\frac{t^2}{n}}\right)^{\frac{n-1}{2}}\int_{0}^{\infty}{u^{\frac{n-1}{2}}e^{-u}\frac{2}{1+\frac{t^2}{n}} du}\\
					&=\frac{1}{\sqrt{2n\pi}2^{\frac{n}{2}}\Gamma{\left(\frac{n}{2}\right)}}\left(\frac{2}{1+\frac{t^2}{n}}\right)^{\frac{n+1}{2}}\int_{0}^{\infty}{u^{\frac{n+1}{2}-1}e^{-u} du} = \frac{1}{\sqrt{2n\pi}2^{\frac{n}{2}}\Gamma{\left(\frac{n}{2}\right)}}\left(\frac{2}{1+\frac{t^2}{n}}\right)^{\frac{n+1}{2}}\Gamma{\left(\frac{n+1}{2}\right)}\\
					&= \frac{\Gamma{\left(\frac{n+1}{2}\right)}}{\sqrt{n\pi}\Gamma{\left(\frac{n}{2}\right)}}\left(1+\frac{t^2}{n}\right)^{-\frac{n+1}{2}} = \text{pdf of} \quad t_n
				\end{align*}
			\end{proof}
		\section{\thesubsection \texorpdfstring{$t$}{t}分布の期待値}
			\begin{shadebox}
				$t$分布は自由度$n \geq 2$の場合に限り期待値が存在し$E[T] = 0$である。
			\end{shadebox}
			\begin{proof}
				\quad\par
				$tf(t)$が$t \to \infty$で$\frac{1}{t}$オーダーより小さくなる条件が$n \geq 2$であることは式を見れば容易にわかる。
				この時に限り$\int_{0}^{\infty}{tf(t) dt}$および$\int_{-\infty}^{0}{tf(t) dt}$が存在し$E[T] = \int_{-\infty}^{\infty}{tf(t) dt}$が存在する。
				被積分関数が奇関数であることから積分値は明らかに0である。
			\end{proof}
			\section{\thesubsection \texorpdfstring{$t$}{t}分布の分散}
			\begin{shadebox}
				$t$分布の分散は自由度$n \leq 2$の場合$\infty$、自由度$n \geq 3$の場合$\frac{n}{n-2}$
			\end{shadebox}
			\begin{proof}
				\quad\par
				\setcounter{equation}{0}
				$n \leq 2$の場合、$t \to\pm\infty$で$t^2f(t)$が$\frac{1}{t}$オーダーを下回らないから$\int_{-\infty}^{\infty}{t^2f(t) dt} = \infty$ \\
				$n \geq 3$の場合、積分は収束する。このとき分散が$\frac{n}{n-2}$であることを示す。\\
				確率密度関数は偶関数だから
				\begin{equation}
					\int_{-\infty}^{\infty}{t^2f(t) dt} = 2\int_{0}^{\infty}{t^2f(t) dt} = \frac{2\Gamma{\left(\frac{n+1}{2}\right)}}{\sqrt{n\pi}\Gamma{\left(\frac{n}{2}\right)}}\int_{0}^{\infty}{t^2\left(1+\frac{t^2}{n}\right)^{-\frac{n+1}{2}} dt}
				\end{equation}
				先に積分の値を計算する。$t^2 = u$なる変数変換を行って
				\[\int_{0}^{\infty}{t^2\left(1+\frac{t^2}{n}\right)^{-\frac{n+1}{2}} dt} = \int_{0}^{\infty}{u\left(1+\frac{u}{n}\right)^{-\frac{n+1}{2}} \frac{1}{2\sqrt{u}}du} = \frac{1}{2}\int_{0}^{\infty}{u^{\frac{1}{2}}\left(1+\frac{u}{n}\right)^{-\frac{n+1}{2}} du}\]
				次に$\displaystyle v = \frac{\frac{u}{n}}{1+\frac{u}{n}} = \frac{u}{n+u}$なる変数変換を行う。$\displaystyle u = \frac{nv}{1-v}$となって$du = \frac{n}{(1-v)^2}dv$となる。積分範囲は$v=0 \to \textcolor{red}{1}$になる。\\
				(※(1)で積分の左側に括りだした定数群にガンマ関数が居る。最終的にこれらが消えてほしいわけだから積分結果がガンマ関数あるいはベータ関数にならなければならない。そこでベータ関数の出現を期待して積分範囲を$0\to\infty$から$0\to 1$に変更するような変数変換を行ったのである。)\\
				上式は
				\begin{align*}
					& \frac{1}{2}\int_{0}^{1}{\left(\frac{nv}{1-v}\right)^{\frac{1}{2}}\left(1+\frac{v}{1-v}\right)^{-\frac{n+1}{2}}\frac{n}{(1-v)^2} dv} = \frac{n\sqrt{n}}{2}\int_{0}^{1}{\sqrt{v}(1-v)^{\frac{n}{2}-2} dv}\\
					&= \frac{n\sqrt{n}}{2}\int_{0}^{1}{v^{\frac{3}{2}-1}(1-v)^{\frac{n-2}{2}-1} dv} = \frac{n\sqrt{n}}{2}B\left(\frac{3}{2},\frac{n-2}{2}\right) = \frac{n\sqrt{n}}{2}\frac{\Gamma{\left(\frac{3}{2}\right)}\Gamma{\left(\frac{n-2}{2}\right)}}{\Gamma{\left(\frac{n+1}{2}\right)}}
				\end{align*}
				ここで
				\[\Gamma{\left(\frac{3}{2}\right)} = \Gamma{\left(\frac{1}{2}+1\right)} = \frac{1}{2}\Gamma{\left(\frac{1}{2}\right)} = \frac{\sqrt{\pi}}{2}\]
				および
				\[\Gamma{\left(\frac{n-2}{2}\right)} = \Gamma{\left(\frac{n}{2}-1\right)}\]
				であり
				\[\Gamma{\left(\frac{n}{2}\right)} = \Gamma{\left(\frac{n}{2}-1+1\right)} = \left(\frac{n}{2}-1\right)\Gamma{\left(\frac{n}{2}-1\right)}\]
				であるから
				\[\Gamma{\left(\frac{n}{2}-1\right)} = \frac{2\Gamma{\left(\frac{n}{2}\right)}}{n-2}\]
				以上の結果を(1)に適用して
				\[\int_{-\infty}^{\infty}{t^2f(t) dt} = \frac{n}{n-2}\]
				最後に
				\[V[T] = E[T^2] - {E[T]}^2 = \int_{-\infty}^{\infty}{t^2f(t) dt} - 0 = \frac{n}{n-2}\]
			\end{proof}
	\chapter{\texorpdfstring{$F$}{F}分布}
		\section{定義}
			\begin{shadebox}
				関数$f(f)$を次のように定義する。
				\[f(f) = \frac{\Gamma{\left(\frac{\nu_1+\nu_2}{2}\right)}}{\Gamma{\left(\frac{\nu_1}{2}\right)}\Gamma{\left(\frac{\nu_2}{2}\right)}}\left(\frac{\nu_1}{\nu_2}\right)^{\frac{\nu_1}{2}}f^{\frac{\nu_1-2}{2}}\left(1+\frac{\nu_1}{\nu_2}f\right)^{-\frac{\nu_1+\nu_2}{2}}\]
				\[\text{where} \hspace{10pt} \nu_1,\nu_2 \in \mathbb{N},\;\;\; f > 0\]
				確率密度関数が上式で与えられるような分布を「自由度$(\nu_1,\nu_2)$の$F$分布」と呼び、$F_{\nu_2}^{\nu_1}$と書くことにする。
			\end{shadebox}
		\section{\texorpdfstring{$X_1,X_2 : \text{indep},\;\;\;X_1 \sim {\chi^2}_{\nu_1},\;\;\;X_2 \sim {\chi^2}_{\nu_2} \hspace{10pt} \Rightarrow \hspace{10pt} F\coloneqq \frac{\frac{X_1}{\nu_1}}{\frac{X_2}{\nu_2}} \sim F_{\nu_2}^{\nu_1}$}{X1,X2: 独立, X1～χ\string^2\_ν1, X2～χ\string^2\_ν2 ⇒ F := (X1/ν1)/(X2/ν2) ～ F\_ν2\string^ν1}}
			\begin{shadebox}
				\[X_1,X_2 : \text{indep},\;\;\;X_1 \sim {\chi^2}_{\nu_1},\;\;\;X_2 \sim {\chi^2}_{\nu_2} \hspace{10pt} \Rightarrow \hspace{10pt} F\coloneqq \frac{\frac{X_1}{\nu_1}}{\frac{X_2}{\nu_2}} \sim F_{\nu_2}^{\nu_1}\]
			\end{shadebox}
			\begin{proof}
				\quad\par
				$X_1,X_2,F$の確率密度関数を$f^{X_1}(x_1),f^{X_2}(x_2),f(f)$と書くことにする。
				\begin{align*}
					P(0 \leq F \leq f) &= P\left(0 \leq X_1 \leq \frac{\nu_1}{\nu_2}x_2f\right) = \int_{x_2=0}^{\infty}{\int_{x_1=0}^{\frac{\nu_1}{\nu_2}x_2f}{f^{X_1}(x_1)f^{X_2}(x_2) dx_1} dx_2}\\
					&= \int_{x_2=0}^{\infty}{f^{X_2}(x_2)\left(\int_{x_1=0}^{\frac{\nu_1}{\nu_2}x_2f}{f^{X_1}(x_1) dx_2}\right) dx_1}
				\end{align*}
				\begin{align}
					\therefore \;\;\; f(f) &= \frac{\partial}{\partial f}P(0 \leq F \leq f) = \int_{x_2=0}^{\infty}{f^{X_2}(x_2)f^{X_1}\left(\frac{\nu_1}{\nu_2}x_2f\right)\frac{\nu_1}{\nu_2}x_2 dx_2} \quad \text{(微分と積分の順序を交換した)} \nonumber \\
					&= \frac{1}{2^{\frac{\nu_2}{2}}\Gamma{\left(\frac{\nu_2}{2}\right)}2^{\frac{\nu_1}{2}}\Gamma{\left(\frac{\nu_1}{2}\right)}}\frac{\nu_1}{\nu_2}\int_{0}^{\infty}{x^{\frac{\nu_2}{2}-1}e^{-\frac{x}{2}}\left(\frac{\nu_1}{\nu_2}f\right)^{\frac{\nu_1}{2}-1}x^{\frac{\nu_1}{2}-1}\exp\left[-\frac{x}{2}\left(1+\frac{\nu_1}{\nu_2}f\right)\right] dx} \nonumber \\
					&=\frac{1}{2^{\frac{\nu_1+\nu_2}{2}}\Gamma{\left(\frac{\nu_2}{2}\right)}\Gamma{\left(\frac{\nu_1}{2}\right)}}\left(\frac{\nu_1}{\nu_2}\right)^\frac{\nu_1}{2}f^{\frac{\nu_1-2}{2}}\int_{0}^{\infty}{x^{\frac{\nu_1+\nu_2}{2}-1}\exp\left[-\frac{x}{2}\left(1+\frac{\nu_1}{\nu_2}f\right)\right] dx}
				\end{align}
				先に積分の値を計算する。$\displaystyle \frac{x}{2}\left(1+\frac{\nu_1}{\nu_2}f\right) = u$なる変数変換を行って
				\begin{align*}
					& \int_{0}^{\infty}{x^{\frac{\nu_1+\nu_2}{2}-1}\exp\left[-\frac{x}{2}\left(1+\frac{\nu_1}{\nu_2}f\right)\right] dx} = \int_{0}^{\infty}{\left(\frac{2}{1+\frac{\nu_1}{\nu_2}f}u\right)^{\frac{\nu_1+\nu_2}{2}-1}e^{-u}\frac{2}{1+\frac{\nu_1}{\nu_2}f} du}\\
					&= \left(1+\frac{\nu_1}{\nu_2}f\right)^{-\frac{\nu_1+\nu_2}{2}}2^{\frac{\nu_1+\nu_2}{2}}\int_{0}^{\infty}{u^{\frac{\nu_1+\nu_2}{2}-1}e^{-u} du} = \left(1+\frac{\nu_1}{\nu_2}f\right)^{-\frac{\nu_1+\nu_2}{2}}2^{\frac{\nu_1+\nu_2}{2}}\Gamma{\left(\frac{\nu_1+\nu_2}{2}\right)}
				\end{align*}
				これを(2)に適用して
				\[f(f) = \frac{\Gamma{\left(\frac{\nu_1+\nu_2}{2}\right)}}{\Gamma{\left(\frac{\nu_1}{2}\right)}\Gamma{\left(\frac{\nu_2}{2}\right)}}\left(\frac{\nu_1}{\nu_2}\right)^{\frac{\nu_1}{2}}f^{\frac{\nu_1-2}{2}}\left(1+\frac{\nu_1}{\nu_2}f\right)^{-\frac{\nu_1+\nu_2}{2}} = \text{pdf of} \quad F_{\nu_2}^{\nu_1}\]
			\end{proof}
		\section{自由度\texorpdfstring{$(\nu_1,\nu_2)$}{(ν1,ν2)}の\texorpdfstring{$F$}{F}分布\texorpdfstring{$F_{\nu_2}^{\nu_1}$}{F\_ν2\string^ν1}は\texorpdfstring{$\nu_2 \geq 3$}{ν2≧3}の時に限り期待値が定義できてその値は\texorpdfstring{$\frac{\nu_2}{\nu_2-2}$}{ν2/(ν2-2)}である。}
			\begin{shadebox}
				自由度$(\nu_1,\nu_2)$の$F$分布$F_{\nu_2}^{\nu_1}$は$\nu_2 \geq 3$の時に限り期待値が定義できてその値は$\frac{\nu_2}{\nu_2-2}$である。
			\end{shadebox}
			\begin{proof}
				\setcounter{equation}{0}
				\quad\par
				$\nu_2 \leq 2$の場合、$x\to \infty$で$xf(x)$が$\frac{1}{x}$オーダーを下回らないので$\int_{0}^{\infty}{xf(x) dx}$は発散し期待値は定義できない。
				$\nu_2 \geq 3$の場合、$x\to \infty$で$xf(x)$が$\frac{1}{x}$オーダーを下回るので積分は収束して期待値が定義できる。この値が$\frac{\nu_2}{\nu_2-2}$であることを示す。\\
				\begin{equation}
					E[F] = \int_{0}^{\infty}{xf(x) dx} = \frac{\Gamma{\left(\frac{\nu_1+\nu_2}{2}\right)}}{\Gamma{\left(\frac{\nu_1}{2}\right)}\Gamma{\left(\frac{\nu_2}{2}\right)}}\left(\frac{\nu_1}{\nu_2}\right)^{\frac{\nu_1}{2}}\int_{0}^{\infty}{x^{\frac{\nu_2}{2}}\left(1+\frac{\nu_1}{\nu_2}x\right)^{-\frac{\nu_1+\nu_2}{2}} dx}
				\end{equation}
				積分の値を先に計算する。$\displaystyle u = \frac{\frac{\nu_1}{\nu_2}x}{1+\frac{\nu_1}{\nu_2}x} = \frac{\nu_1x}{\nu_2+\nu_1x}$なる変数変換を行うと$x = \frac{\nu_2}{\nu_1}\frac{u}{1-u}, \quad dx = \frac{\nu_2}{\nu_1}\frac{1}{(1-u)^2}du$となり積分範囲は$u=0\to \textcolor{red}{1}$になる。\\
				(※(1)で積分の左側に括りだした定数群にガンマ関数が居る。最終的にこれらが消えてほしいわけだから積分結果がガンマ関数あるいはベータ関数にならなければならない。そこでベータ関数の出現を期待して積分範囲を$0\to\infty$から$0\to 1$に変更するような変数変換を行ったのである。)\\
				(1)の積分は
				\begin{align*}
					& \int_{0}^{1}{\left(\frac{\nu_2}{\nu_1}\right)^{\frac{\nu_1}{2}}\left(\frac{u}{1-u}\right)^{\frac{\nu_1}{2}}\left(1+\frac{u}{1-u}\right)^{-\frac{\nu_1+\nu_2}{2}}\frac{\nu_2}{\nu_1}\frac{1}{(1-u)^2} du}\\
					&= \left(\frac{\nu_1}{\nu_2}\right)^{-\frac{\nu_1}{2}-1}\int_{0}^{1}{u^{\frac{\nu_1}{2}}(1-u)^{\frac{\nu_2}{2}-2} du} = \left(\frac{\nu_1}{\nu_2}\right)^{-\frac{\nu_1}{2}-1}\int_{0}^{1}{u^{\frac{\nu_1}{2}+1-1}(1-u)^{\frac{\nu_2}{2}-1-1} du}\\
					&= \left(\frac{\nu_1}{\nu_2}\right)^{-\frac{\nu_1}{2}-1}B\left(\frac{\nu_1}{2}+1,\frac{\nu_2}{2}-1\right) = \left(\frac{\nu_1}{\nu_2}\right)^{-\frac{\nu_1}{2}-1}\frac{\Gamma{\left(\frac{\nu_1}{2}+1\right)}\Gamma{\left(\frac{\nu_2}{2}-1\right)}}{\Gamma{\left(\frac{\nu_1+\nu_2}{2}\right)}}\\
					&= \left(\frac{\nu_1}{\nu_2}\right)^{-\frac{\nu_1}{2}-1}\frac{\frac{\nu_1}{2}\Gamma{\left(\frac{\nu_1}{2}\right)}\Gamma{\left(\frac{\nu_2}{2}\right)}}{\left(\frac{\nu_2}{2}-1\right)\Gamma{\left(\frac{\nu_1+\nu_2}{2}\right)}}
				\end{align*}
				これを(1)に適用して
				\[E[F] = \frac{\nu_2}{\nu_2-2}\]
			\end{proof}
		\section{自由度\texorpdfstring{$(\nu_1,\nu_2)$}{(ν1,ν2)}の\texorpdfstring{$F$}{F}分布\texorpdfstring{$F_{\nu_2}^{\nu_1}$}{F\_ν2\string^ν1}の分散は\texorpdfstring{$\nu_2 \leq 4$}{ν2<4}の場合\texorpdfstring{$\infty$}{∞}、\texorpdfstring{$\nu_2 \geq 5$}{ν2≧5}の場合\texorpdfstring{$\displaystyle \frac{2{\nu_2}^2(\nu_1+\nu_2-2)}{\nu_1(\nu_2-2)^2(\nu_2-4)}$}{(2(ν2)\string^2 (ν1+ν2-2))/(ν1(ν2-2)\string^2 (ν2-4))}である。}
			\begin{shadebox}
				自由度$(\nu_1,\nu_2)$の$F$分布$F_{\nu_2}^{\nu_1}$の分散は$\nu_2 \leq 4$の場合$\infty$、$\nu_2 \geq 5$の場合$\displaystyle \frac{2{\nu_2}^2(\nu_1+\nu_2-2)}{\nu_1(\nu_2-2)^2(\nu_2-4)}$である。
			\end{shadebox}
			\begin{proof}
				\setcounter{equation}{0}
				\quad\par
				$\nu_2 \leq 4$の場合、$x\to \infty$で$x^2f(x)$が$\frac{1}{x}$オーダーを下回らないので$\int_{0}^{\infty}{x^2f(x) dx} = \infty$\\
				$\nu_2 \geq 5$の場合、$x\to \infty$で$x^2f(x)$が$\frac{1}{x}$オーダーを下回るので積分は収束する。このとき分散が$\displaystyle \frac{2{\nu_2}^2(\nu_1+\nu_2-2)}{\nu_1(\nu_2-2)^2(\nu_2-4)}$であることを示す。\\
				$V[F] = E[F^2] - {E[F]}^2$であり第2項は定理2より既知である。第1項を計算する。
				\begin{equation}
					E[F^2] = \int_{0}^{\infty}{x^2f(x) dx} = \frac{\Gamma{\left(\frac{\nu_1+\nu_2}{2}\right)}}{\Gamma{\left(\frac{\nu_1}{2}\right)}\Gamma{\left(\frac{\nu_2}{2}\right)}}\left(\frac{\nu_1}{\nu_2}\right)^{\frac{\nu_1}{2}}\int_{0}^{\infty}{x^{\frac{\nu_2}{2}+1}\left(1+\frac{\nu_1}{\nu_2}x\right)^{-\frac{\nu_1+\nu_2}{2}} dx}
				\end{equation}
				積分の値を先に計算する。$\displaystyle u = \frac{\frac{\nu_1}{\nu_2}x}{1+\frac{\nu_1}{\nu_2}x} = \frac{\nu_1x}{\nu_2+\nu_1x}$なる変数変換を行うと$x = \frac{\nu_2}{\nu_1}\frac{u}{1-u}, \quad dx = \frac{\nu_2}{\nu_1}\frac{1}{(1-u)^2}du$となり積分範囲は$u=0\to \textcolor{red}{1}$になる。\\
				(1)の積分は
				\begin{align*}
					& \int_{0}^{1}{\left(\frac{\nu_2}{\nu_1}\right)^{\frac{\nu_1}{2}+1}\left(\frac{u}{1-u}\right)^{\frac{\nu_1}{2}+1}\left(1+\frac{u}{1-u}\right)^{-\frac{\nu_1+\nu_2}{2}}\frac{\nu_2}{\nu_1}\frac{1}{(1-u)^2} du}\\
					&= \left(\frac{\nu_1}{\nu_2}\right)^{-\frac{\nu_1}{2}-2}\int_{0}^{1}{u^{\frac{\nu_1}{2}+1}(1-u)^{\frac{\nu_2}{2}-3} du} = \left(\frac{\nu_1}{\nu_2}\right)^{-\frac{\nu_1}{2}-2}B{\left(\frac{\nu_1}{2}+2,\frac{\nu_2}{2}-2\right)}\\
					&= \left(\frac{\nu_1}{\nu_2}\right)^{-\frac{\nu_1}{2}-2}\frac{\Gamma{\left(\frac{\nu_1}{2}+2\right)}\Gamma{\left(\frac{\nu_2}{2}-2\right)}}{\Gamma{\left(\frac{\nu_1+\nu_2}{2}\right)}}\\
					&= \left(\frac{\nu_1}{\nu_2}\right)^{-\frac{\nu_1}{2}-2}\frac{\left(\frac{\nu_1}{2}+1\right)\frac{\nu_1}{2}\Gamma{\left(\frac{\nu_1}{2}\right)}}{\Gamma{\left(\frac{\nu_1+\nu_2}{2}\right)}}\frac{4\Gamma{\left(\frac{\nu_2}{2}\right)}}{(\nu_2-4)(\nu_2-2)}
				\end{align*}
				これを(1)に適用して
				\[E[F^2] = \frac{{\nu_2}^2(\nu_1+2)}{\nu_1(\nu_2-4)(\nu_2-2)}\]
				定理2で示された$E[F]$を用いて
				\[V[F] = E[F^2]-{E[F]}^2 = \frac{{\nu_2}^2(\nu_1+2)}{\nu_1(\nu_2-4)(\nu_2-2)} - \frac{{\nu_2}^2}{(\nu_2-2)^2} = \frac{2{\nu_2}^2(\nu_1+\nu_2-2)}{\nu_1(\nu_2-2)^2(\nu_2-4)}\]
			\end{proof}
	\chapter{推定}
		\section{標本分散}
			\subsection{\texorpdfstring{$\E{}{\frac{1}{n}\sum_{i=1}^{n}{(X_i-\overline{X})^2}} = \frac{n-1}{n}\sigma^2$}{E((1/n)Σ\_\{i=1\}\string^n (Xi - sampleAvg(X))\string^2) = ((n-1)/n)σ\string^2}}
				\begin{shadebox}
					$X$が平均$\mu$,分散$\sigma^2$の母集団に属しているとき、そこからランダムに抽出した標本$(X_1,X_2,\dots,X_n)$の標本分散$\displaystyle S^2 \coloneqq \frac{1}{n}\sum_{i=1}^{n}{(X_i-\overline{X})^2}$の期待値は$\displaystyle \text{E}(S^2) = \frac{n-1}{n}\sigma^2$である。
				\end{shadebox}
				\begin{proof}
				{%スコープ。\s 上書き回避
					\newcommand{\EE}[1]{\text{E}^2\left[#1\right]}
					\renewcommand{\s}{\sigma^2}
					\newcommand{\mX}{\overline{X}}
					\begin{equation}
						\E{}{S^2} = \E{}{\frac{1}{n}\sum_{i=1}^{n}{(X_i-\mX)^2}} = \frac{1}{n}\sum_{i=1}^{n}{\E{}{(X_i-\mX)^2}} = \frac{1}{n}\sum_{i=1}^{n}{\left(\E{}{X_i^2}+\E{}{\mX^2}-2\E{}{\mX X_i}\right)}
					\end{equation}
					然るに\\
					$\Var{}{X_i} = \E{}{X_i^2}-\EE{X_i}$より
					\begin{equation}
						\E{}{X_i^2} = \Var{}{X_i}+\EE{X_i} = \s+\mu^2
					\end{equation}
					\begin{align}
						\E{}{\mX^2} &= \frac{1}{n^2}\E{}{\left(\sum_{i=1}^{n}{X_i}\right)^2} = \frac{1}{n^2}\E{}{\sum_{i=1}^{n}{X_i^2} + \sum_{i\neq j}^{}{X_iX_j}} = \frac{1}{n^2}\sum_{i=1}^{n}{\E{}{X_i^2}} + \frac{1}{n^2}\sum_{i\neq j}^{}{\E{}{X_iX_j}} \nonumber \\
						&= \frac{1}{n^2}\sum_{i=1}^{n}{(\s+\mu^2)} + \frac{1}{n^2}\sum_{i\neq j}^{}{\E{}{X_i}\E{}{X_j}} = \frac{1}{n^2}\sum_{i=1}^{n}{(\s+\mu^2)} + \frac{1}{n^2}\sum_{i\neq j}^{}{\mu^2} \nonumber \\
						&= \frac{1}{n}(\s+\mu^2) + \frac{1}{n^2}(n^2-n)\mu^2 = \frac{1}{n}\s + \mu^2
					\end{align}
					\begin{align}
						\E{}{\mX X_i} &= \frac{1}{n}\E{}{(X_1+X_2+\dotsb+X_n)X_i} = \frac{1}{n}\E{}{X_i^2 + \sum_{j=1,\neq i}^{n}{X_jX_i}} \nonumber \\
						&= \frac{1}{n}\E{}{X_i^2} + \frac{1}{n}\sum_{j=1,\neq i}^{n}{\E{}{X_jX_i}} = \frac{1}{n}(\s+\mu^2) + \frac{1}{n}\sum_{j=1,\neq i}^{n}{\E{}{X_j}\E{}{X_i}} \nonumber \\
						&= \frac{1}{n}(\s+\mu^2) + \frac{1}{n}\sum_{j=1,\neq i}^{n}{\mu^2} = \frac{1}{n}(\s+\mu^2) + \frac{1}{n}(n-1)\mu^2 = \frac{1}{n}\s + \mu^2
					\end{align}
					(2),(3),(4)を(1)に適用して
					\[\E{}{\s} = \frac{1}{n}\sum_{i=1}^{n}{\left(\s-\frac{1}{n}\s\right)} = \frac{n-1}{n}\s\]
				}
				\end{proof}
